\documentclass{article} % For LaTeX2e
\usepackage{listings}
\usepackage[svgnames,dvipsnames]{xcolor} % Specify colors by their 'svgnames', for a full list of all colors available see here: http://www.latextemplates.com/svgnames-colors
\usepackage{etoolbox}
\usepackage{nips15submit_e,times}
\usepackage{hyperref}
\usepackage{url}
\usepackage{tcolorbox}
\usepackage{tikz}
\usepackage{pgfplots}
\usepackage{bm}
\usepackage{amsmath,amssymb}
\usepackage{natbib}
%\pgfplotsset{compat=1.9}
\usepackage[framemethod=TikZ]{mdframed}% http://ctan.org/pkg/mdframed
\usepackage{graphicx}
\usepackage{caption}
\usepackage{subcaption}

\usepackage{fancyvrb}
\usepackage{placeins}
\usetikzlibrary{tikzmark}
\usetikzlibrary{bayesnet}


\newcommand{\argmax}[1]{\underset{#1}{\operatorname{arg}\,\operatorname{max}}\;}


\apptocmd\normalsize{%
 \abovedisplayskip=5pt plus 3pt minus 3pt
 \abovedisplayshortskip=0pt 
 \belowdisplayskip=5pt plus 3pt minus 3pt
 \belowdisplayshortskip=0pt
}{}{}

\usetikzlibrary{pgfplots.groupplots}
%opening
%\documentstyle[nips14submit_09,times,art10]{article} % For LaTeX 2.09

\newcommand{\gpmem}{\texttt{gpmem}}
\newcommand{\emu}{{\textrm{emu}}}
\newcommand{\restr}{{\textrm{restr}}}
\newcommand{\true}{{\textrm{true}}}
\newcommand{\rmnew}{{\textrm{new}}}
\newcommand{\past}{{\textrm{past}}}
\newcommand{\prior}{{\textrm{prior}}}
\newcommand{\noisy}{{\textrm{noisy}}}
\newcommand{\noise}{{\textrm{noise}}}

\newcommand{\Acal}{\mathcal{A}}
\newcommand{\R}{\mathbb{R}}

\newcommand{\abf}{\mathbf{a}}
\newcommand{\fbf}{\mathbf{f}}
\newcommand{\rbf}{\mathbf{r}}
\newcommand{\wbf}{\mathbf{w}}
\newcommand{\xbf}{\mathbf{x}}
\newcommand{\ybf}{\mathbf{y}}
\newcommand{\Kbf}{\mathbf{K}}
\newcommand{\Ibf}{\mathbf{I}}

\newcommand{\pn}[1]{\left( #1 \right)}
\newcommand{\bkt}[1]{\left[ #1 \right]}
\newcommand{\br}[1]{\left\{ #1 \right\}}
\newcommand{\abs}[1]{\left\lvert #1 \right\rvert}
\newcommand{\Ebkt}[2][]{\mathbb{E}_{#1}\bkt{#2}}
\newcommand{\mvert}{\ \middle\vert\ }

\newcommand{\paperOrChapter}{paper}

\DeclareMathOperator*{\Cov}{Cov}

% The following is a modified version of what's available on Wikibooks,
% http://en.wikibooks.org/wiki/LaTeX/Source_Code_Listings
  \definecolor{mygreen}{rgb}{0,0.6,0}
  \definecolor{mygray}{rgb}{0.5,0.5,0.5}
  \definecolor{mymauve}{rgb}{0.58,0,0.82}

  \lstset{
    language=LISP,
    comment=[l]{//},
    keywordstyle=\color{black},
    commentstyle=\color{mygreen},
    keywordstyle=[2]\color{RawSienna},
    keywords=[2]{assume,predict,infer,observe,define},
    keywordstyle=[3]\color{black},
    keywords=[3]{then,else,proc,tag,gamma,uniform_continuous,log,make_squaredexp,make_whitenoise,gpmem,mapv,first,repeat,do,mh,one,add_funcs,make_se,lambda,lookup,argmax_of_array,linspace,run,array,pass,mc_argmax,uniform_structure,subset,lte,if,flip,mult_funcs,for,to,get_neal_blackbox,get_neal_data_xs,get_data_xs,size,get_bayesopt_blackbox},
    basicstyle=\tt,
    literate=%
    *{0}{{{\color{DarkBlue}0}}}1
    {1}{{{\color{DarkBlue}1}}}1
    {2}{{{\color{DarkBlue}2}}}1
    {3}{{{\color{DarkBlue}3}}}1
    {4}{{{\color{DarkBlue}4}}}1
    {5}{{{\color{DarkBlue}5}}}1
    {6}{{{\color{DarkBlue}6}}}1
    {7}{{{\color{DarkBlue}7}}}1
    {8}{{{\color{DarkBlue}8}}}1
    {9}{{{\color{DarkBlue}9}}}1
  }
\title{Probabilistic Programming with Gaussian Process Memoization}


\author{
}

% The \author macro works with any number of authors. There are two commands
% used to separate the names and addresses of multiple authors: \And and \AND.
%
% Using \And between authors leaves it to \LaTeX{} to determine where to break
% the lines. Using \AND forces a linebreak at that point. So, if \LaTeX{}
% puts 3 of 4 authors names on the first line, and the last on the second
% line, try using \AND instead of \And before the third author name.

\newcommand{\fix}{\marginpar{FIX}}
\newcommand{\new}{\marginpar{NEW}}

%\nipsfinalcopy % Uncomment for camera-ready version

\begin{document}


\maketitle

\begin{abstract}
This paper describes the {\em Gaussian process memoizer}, a probabilistic programming technique that uses Gaussian processes to provide a statistical alternative to memorization. Memoizing a target procedure results in a “self-caching” wrapper that remembers previously computed values. Gaussian process memoization additionally produces a statistical emulator based on a Gaussian process whose predictions automatically improve whenever a new value of the target procedure becomes available. This paper also introduces  an efficient implementation, named {\tt gpmem}, that can use kernels given by a broad class of probabilistic programs. The flexibility of {\tt gpmem} is illustrated via three applications: (i) GP regression with hierarchical hyper-parameter learning, (ii) Bayesian structure learning via compositional kernels generated by a probabilistic grammar, and (iii) a bandit formulation of Bayesian optimization with automatic inference and action selection. All applications share a single 50-line Python library and require fewer than 20 lines of probabilistic code each.
\end{abstract}

\section{Introduction}
\ac{GP} are widely used tools in statistics~\citep{barry1986nonparametric}, machine learning~\citep{neal1995bayesian,williams1998bayesian,kuss2005assessing,rasmussen2006gaussian,damianou2013deep}, robotics \citep{ferris2006gaussian}, computer vision~\citep{kemmler2013one}, and scientific computation~\citep{kennedy2001bayesian,schneider2008simulations,kwan2013cosmic}.
% ToDo: get a citation of neal on stats!
They are also central to probabilistic numerics, an emerging effort to develop more computationally efficient numerical procedures, and to Bayesian optimization, a family of meta-optimization techniques that are widely used to tune parameters for deep learning algorithms~\citep{snoek2012practical,gelbart2014bayesian}. They have even seen use in artificial intelligence; for example, they provide the key technology behind a project that produces qualitative natural language descriptions of time series~\citep{duvenaud2013structure,lloyd2014automatic}.

This paper describes Gaussian process memoization, a technique for integrating \ac{GP}s into a probabilistic programming language, and demonstrates its utility by re-implementing and extending state-of-the-art applications of the GP. Memoization is a classic programming technique in which a procedure is augmented with an input-output cache that is checked each time before the function is invoked. This prevents unnecessary recomputation, potentially saving time at the cost of increased storage requirements. Gaussian process memoization, implemented by the {\tt gpmem()} procedure, generalizes this idea to include a statistical emulator that uses previously computed values as data in a statistical model that can accurately forecast probable outputs. The covariance function for the Gaussian process is also allowed to be an arbitrary probabilistic program.

This paper presents three applications of gpmem: (i) a replication of results~\citet{neal1997monte} on outlier rejection via hyper-parameter inference; (ii) a fully Bayesian extension to the Automated Statistician project; and (iii) an implementation of Bayesian optimization via Thompson sampling. The first application can in principle be replicated in several other probabilistic languages embedding the proposal that is described in this paper. The remaining two applications rely on distinctive capabilities of Venture~\citep{mansinghka2014venture}: support for fully Bayesian structure learning and language constructs for inference programming. All applications share a single 50-line Python library and require fewer than 20 lines of probabilistic code each.
% ToDo: the 50 lines are bit understated

\subsection{Related Work}
A recent review on artificial intelligence and machine learning in Natur named
\ac{GP}s and probabilistic programming as amongst the most promising directions for 
future research \citep{ghahramani2015probabilistic}. A practitioner can find many software packages
for \ac{GP}s available for free~\citep[e.g.][]{rasmussen2010gaussian,gpy2014,PyGPs}.
However, the context that a practitioner can provide with them is mostly limited to regression and
classification tasks. Yet, \ac{GP}s performed well in applications that go much
further than these tasks. Neal shows an interesting application of Bayesian regression in
the presence of outliers with a hierarchical system of hyper-priors~\citet{neal1997monte}.
\ac{GP}s have been applied to emulated complex computation such as in the case of
\ac{CCF}~\citep{kwan2013cosmic}. At the heart of this project is the
\ac{CCF} that exploits accuracy of large scale  high-resolution cosmological
simulations that are computationally expensive. Here, a sophisticated sampling scheme
provides an optimal sampling strategy for the cosmological models to be
simulated. The measurements from the simulations are then translated into functions that can
be easily interpolated with \ac{GP}s.

The discovery of symbolic relations and concepts has also been the
focus of attention of Bayesian non-parametrics~\citep[e.g.][]{kemp2006learning}.
The Automatic Statistician introduces inductive learning of
symbolic expressions by kernel structure learning, through performing a greedy search over
the space of possible \ac{GP} kernel
structure~\citep{duvenaud2013structure,lloyd2014automatic}. The resulting
structure comes with symbolic interpretation that can be understood by humans. The
engineering involved a significant work-around for the
GPML-toolbox~\citep{rasmussen2010gaussian} which, whilst providing a novel
solution, still struggles to fulfill requirements to be fully Bayesian.

A further example of advanced use of \ac{GP}s beyond regression and
classification is the seemingly deterministic domain of
optimization. To be Bayesian over the values of expensive functions, \ac{GP}s were used
as emulators, leading to efficient optimization of otherwise expensive machine learning algorithms~\citep{snoek2012practical}. Bayesian optimization
updates the belief system of a Bayesian agent when new information about the
true underlying function to optimize becomes available. A sampling scheme for
performing Bayesian Optimization in very high dimensions has been introduced
by~\citep{mahendran2012adaptive}.  


Probabilistic programming systems on the other hand have pushed the boundaries
of models usable for machine learning with general purpose inference machinery.
Examples include picture, which is combining fast data-driven methods for image
interpretation such as deep neural networks with generative
modelling~\citep{kulkarni2015picture}. Classification with large, deep and evolving hierarchical
models as needed for the identification and tracking of resident space objects
was enabled with the Figaro language~\citep{ruttenberg2014hierarchical}.
Nitti and colleagues present a probabilistic programming approach to
planning with \ac{MDP}s in hybrid domains, that is discrete and continuous-valued
domains as well as domains with an unknown number of
objects~\citep{nitti2015planning}. Bayesian Program
induction outperforms deep neural networks in tasks with only one training example provided~\citep{lake2015humanlevel}. This kind of learning is
particularly hard for machine learning while at the same time considerably easy
for humans. 



\section{Background}
\subsection{Memoization}
\input{sections/Memoization.tex}
\subsection{Venture}
Venture is a compositional language for custom inference strategies that comes with a Scheme-like and a JavaScript-like front-end syntaxes.
Its implementation is based on on three concepts:
  (i) \emph{stochastic procedures} that specify and encapsulate random variables, analogously to conditional probability tables in a Bayesian network;
  (ii) \emph{execution traces} that represent (partial) execution histories and track the conditional dependencies of the random variables occurring therein; and
  (iii) \emph{scaffolds} that partition execution histories and factor global inference problems into sub-problems.
These building blocks provide a powerful and concise way to represent probability distributions, including distributions with a dynamically determined and unbounded set of random variables.
In this \paperOrChapter we will use only the four basic Venture directives: ASSUME, OBSERVE, SAMPLE and INFER.
\begin{itemize}
  \item ASSUME induces a hypothesis space for (probabilistic) models including random variables by binding the result of a supplied expression to a supplied symbol.
  \item Whereas in Scheme an expression is evaluated within an environment, in Venture an expression is evaluated within a (partial) trace of the model program.
    Thus, the value of an expression within a model program is a random variable, whose randomness comes from the distribution on possible execution traces of the program.
    The SAMPLE directive samples the value of the supplied expression within the current model program.
  \item OBSERVE constrains the supplied expression to have the supplied value.
    In other words, all samples taken after an OBSERVE are conditioned on the observed data.
  \item INFER uses the supplied inference program to mutate the execution trace.
    For a correct inference program, this will result approximate sampling from the true posterior on execution traces, conditioned on the model and constraints introduced by ASSUME and OBSERVE.
    The posterior on any random variable can then be approximately sampled by calling SAMPLE to extract values from the trace.
\end{itemize}

INFER is commonly done using the Metropolis--Hastings algorithm (MH)~\citep{metropolis1953equation}.
Many of the most popular MCMC algorithms can be interpreted as special cases of MH~\citep{andrieu2003introduction}.
We can outline the MH algorithm as follows.
The following two-step process is repeated as long as desired (say, for $T$ iterations):
First we sample $x^*$ from a proposal distribution $q$:
\begin{equation}
 x^* \sim q(x^* \mid x^{t});
\end{equation}
then we accept this proposal ($x^{t+1} \leftarrow x ^*$) with probability
\begin{equation}
  \alpha = \min \br{
    1,
    \frac{p(x^*) q(x^{t}\mid x^*)}{p(x^{t}) q(x^* \mid x^{t})}};
\end{equation}
if the proposal is not accepted then we take $x^{t+1} \gets x^t$.

Venture includes a built-in generic MH inference program which performs the above steps on any specified set of random variables in the model program.
In that inference program, partial execution traces play the role of $x$ above.




\subsection{Gaussian Processes}
\ac{GP}s are a Bayesian method for regression. We consider the regression input to be real-valued scalars $x_i$ and the regression output as the value of a function $f$ at $x_i$. The complete training data will be denoted by column vectors $\mathbf{x}$ and $\mathbf{f}$. Unseen test data is denoted with $\mathbf{\hat{x}}$ and $\mathbf{\hat{f}}$.
\ac{GP}s present a non-parametric way to express prior knowledge on the space of all possible functions $f$ modeling
a regression relationship.
Formally, a GP is an infinite-dimensional extension of the multivariate Gaussian distribution.
\begin{definition}
A Gaussian process (GP) is a collection of random variables, any
finite number of which have a joint Gaussian distribution~\citep[][chapter 2]{rasmussen2006gaussian}.
\end{definition}
The collection of random variables $\br{f(x)}$ (indexed by $x$) represents the
values of the function $f$ at each location $x$.
We write $f \sim \ac{GP}(m,k)$, where $m$ is the {\em mean function} and $k$ is the {\em covariance function} or {\em kernel}.
That is, $m(x)$ is the prior mean of the random variable $f(x)$, and $k(x,x')$ is the prior covariance of the random variables $f(x)$ and $f(x')$.
Throughout, we write $\mathbf{K}(\xbf,\xbf^\prime)$ for the prior covariance matrix determined by $\xbf$ and $\xbf^\prime$, that is, the covariance between the random vectors $\{f(x)\}_{x \in \xbf}$ and $\{f(x')\}_{x' \in \xbf'}$.

We now compute the predictive posterior distribution of test output $\hat\fbf := f(\hat\xbf)$ conditioned on training data $\fbf := f(\xbf)$.  (Here $\xbf$ and $\hat\xbf$ are known constant vectors, and we are conditioning on an observed value of $\fbf$.)  To simplify the calculation, we will assume the prior mean $m$ is identically zero; once the derivation is done, this assumption can be easily relaxed via translation.

The predictive posterior can be computed by first forming the joint density when both training and test data are treated as randomly chosen from the prior, then fixing the value of $\fbf$ to a constant.  To start, let
\[
  \Sigma := \bmat{
    K(\xbf, \xbf)     & K(\xbf, \hat\xbf)     \\
    K(\hat\xbf, \xbf) & K(\hat\xbf, \hat\xbf)
  }
  \text{ and }
  \Sigma^{-1} =: \bmat{
    M_{11} & M_{12} \\
    M_{21} & M_{22}
  }.
\]
We then have
\[
  P(\fbf, \hat\fbf)
  \propto
  \exp\br{
    -\frac12
    \bmat{\fbf^\top & \hat\fbf^\top}
    \bmat{M_{11} & M_{12} \\ M_{21} & M_{22}}
    \bmat{\fbf \\ \hat\fbf}
  }.
\]
Treating $\fbf$ as a fixed constant, we obtain
\[
  P\pn{\hat\fbf \mvert \fbf}
  \propto
  P(\fbf, \hat\fbf)
  \propto
  \exp\br{
    -\frac12 \hat\fbf^\top M_{22} \hat\fbf
    - \hbf^\top \hat\fbf
  },
\]
where $\hbf = M_{21} \fbf$ is a constant vector.  Thus $P(\hat\fbf | \fbf)$ is Gaussian,
\begin{equation}
  P\pn{\hat\fbf \mvert \fbf} \sim \Ncal(\hat\mubf, \hat\Kbf),
\end{equation}
with covariance matrix $\hat\Kbf = M_{22}^{-1}$.  To find its mean $\hat\mubf$, we note that $P_{\hat\fbf|\fbf}(\hat\fbf + \hat\mubf)$ is Gaussian with the same covariance as $P(\hat\fbf | \fbf)$, but its exponent has no linear term:
\begin{align*}
  P_{\hat\fbf|\fbf} \pn{\hat\fbf + \hat\mubf \mvert \fbf}
  &\propto
  \exp\br{
    -\frac12 (\hat\fbf + \hat\mubf)^\top M_{22} (\hat\fbf + \hat\mubf)
    - \hbf^\top (\hat\fbf + \hat\mubf)
  } \\
  &\propto
  \exp\br{
    -\frac12 \hat\fbf^\top M_{22} \hat\fbf
    - \underbrace{(\hbf + M_{22} \hat\mubf)^\top}_{\text{must be $0$}} \hat\fbf
  }.
\end{align*}
Thus $\hbf = -M_{22} \hat\mubf$ and $\hat\mubf = -M_{22}^{-1} \hbf = -M_{22}^{-1} M_{21} \fbf$.

The partioned inverse equations (\citealp*{barnett1979matrix} following \citealp*{mackay1998introduction}) give
\begin{align*}
  M_{22}^{-1} &= K(\hat\xbf,\hat\xbf) - K(\hat\xbf,\xbf) K(\xbf,\xbf)^{-1} K(\xbf,\hat\xbf), \\
  M_{21} &= -M_{22} K(\hat\xbf,\xbf) K(\xbf,\xbf)^{-1}.
\end{align*}
Substituting these in the above gives
\begin{align}
  \hat\Kbf &= K(\hat\xbf,\hat\xbf) - K(\hat\xbf,\xbf) K(\xbf,\xbf)^{-1} K(\xbf,\hat\xbf), \\
  \hat\mubf &= K(\hat\xbf,\xbf) K(\xbf,\xbf)^{-1}\fbf.
\end{align}

Often one assumes the observed regression output is noisily measured, that is, one only sees the values of $\ybf_\noisy = \mathbf{f}+ \wbf$ where $\wbf$ is Gaussian white noise with variance $\sigma_\noise^2$. This noise term can be absorbed into the covariance matrix $\mathbf{K}(\mathbf{x},\mathbf{x})$ which in the following, we will write as $\mathbf{K}$ for readability. The log-likelihood of a \ac{GP} can then be written as:
\begin{equation}
\label{eq:gplogdens}
\log P(\mathbf{f} \mid \xbf) =
-\frac12 \ybf^\top 
\mathbf{K}^{-1} \ybf
- \frac12\log \abs{\mathbf{K}}
- \frac{n}{2}\log 2\pi
\end{equation}
where $n$ is the number of data points.
Both log-likelihood and predictive posterior can be computed efficiently using a \ac{SP} in Venture~\citep{mansinghka2014venture}
with an algorithm that resorts to Cholesky factorization\citep[chap. 2]{rasmussen2006gaussian}. 
We write the Cholesky factorization as 
$\mathbf{L} \coloneqq \text{chol}(\mathbf{K})$ when
:
\begin{equation}
\mathbf{K} = LL^\top
\end{equation}
where L is a lower triangular matrix. This allows us to compute the inverse of a covariance matrix as
\begin{equation}
\mathbf{K}^{-1} = (\mathbf{L}^{-1})^\top (\mathbf{L}^{-1})
\end{equation}
and its determinant as 
\begin{equation}
det(\mathbf{K}) = det(\mathbf{L})^2
\end{equation}
We compute (\ref{eq:gplogdens}) as
\begin{equation}
\log(P(\mathbf{f}\mid \mathbf{x})\coloneqq - \frac{1}{2} \mathbf{f}^\top \bm{\alpha} - \sum_i \log \mathbf{L}_{ii} - \frac{n}{2} \log 2 \pi
\end{equation}
where 
\begin{equation}
\label{eq:chol_L}
\mathbf{L} \coloneqq \text{chol}(\mathbf{K})
\end{equation}
and 
\begin{equation}
\label{eq:alpha}
\bm{\alpha} \coloneqq  \mathbf{L}^\top \backslash(\mathbf{L} \backslash \mathbf{f}). 
\end{equation}
This results in a computational complexity of $\mathcal{O}(n^3)$ in the number of data points for
sampling with a complexity of $n^3/6$ for (\ref{eq:chol_L}) an $n^2/2$ for (\ref{eq:alpha}). 



The covariance function (or kernel) of a \ac{GP} governs high-level properties of the observed data such as linearity, periodicity and smoothness.
It comes with few free parameters that we call hyper-parameters.
Adjusting the hyper-parameters changes non-qualitative attributes such as length
scales while preserving the qualitative properties of the distribution.
These high-level properties are compositional via addition and multiplication of different covariance functions. An addition models a global interaction, that is an interaction of two high-level components that is qualitatively not dependent on the input space. An example for this a periodic function with a linear trend.
A multiplication models a local interaction of two components. 
An example for this is a periodic function with a linearly increasing amplitude. We demonstrate kernel composition with local interactions in the tutorial in Fig. \ref{fig:composition_tutorial}. 


\begin{figure}
 \centering
     \begin{subfigure}[b]{0.3\textwidth}
        \includegraphics[width=\textwidth]{figs/composition/composition_demo_LIN_prior.png}
        \caption{LIN}
    \end{subfigure}
    ~ %add desired spacing between images, e. g. ~, \quad, \qquad, \hfill etc. 
      %(or a blank line to force the subfigure onto a new line)
    \begin{subfigure}[b]{0.3\textwidth}
        \includegraphics[width=\textwidth]{figs/composition/composition_demo_PER_prior.png}
        \caption{PER}
    \end{subfigure}
    ~ %add desired spacing between images, e. g. ~, \quad, \qquad, \hfill etc. 
    %(or a blank line to force the subfigure onto a new line)
    \begin{subfigure}[b]{0.3\textwidth}
        \includegraphics[width=\textwidth]{figs/composition/composition_demo_LINxPER_prior.png}
        \caption{LIN $\times$ PER}

    \end{subfigure}   \begin{subfigure}[b]{0.3\textwidth}
        \includegraphics[width=\textwidth]{figs/composition/composition_demo_LIN.png}
        \caption{LIN}
    \end{subfigure}
    ~ %add desired spacing between images, e. g. ~, \quad, \qquad, \hfill etc. 
      %(or a blank line to force the subfigure onto a new line)
    \begin{subfigure}[b]{0.3\textwidth}
        \includegraphics[width=\textwidth]{figs/composition/composition_demo_PER.png}
        \caption{PER}
    \end{subfigure}
    ~ %add desired spacing between images, e. g. ~, \quad, \qquad, \hfill etc. 
    %(or a blank line to force the subfigure onto a new line)
    \begin{subfigure}[b]{0.3\textwidth}
        \includegraphics[width=\textwidth]{figs/composition/composition_demo_LINxPER.png}
        \caption{LIN $\times$ PER}
    \end{subfigure}

%20.0739791735
%6.31647597198
%37.7184218042
%19.1051376016

\caption{We depict kernel composition. 
(a) shows raw data (black) generated with a sine function with linearly growing amplitude (blue). 
(b) shows the linear and the periodic base kernel as well as a composition of both. 
The multiplication of the two kernels indicates local interaction. The local interaction we account for in this case is the growing amplitude (a). (c-e) show the parameterized kernels that we introduced in (b).
The parameters are sampled from the posterior distribution on parameters for a \ac{GP} with kernel $\mathbf{K}=\text{LIN} \times \text{PER}$ where the data points from (a) are observed.
(f-h) show samples from the prior where the parametrization from (c-e) is used, that is, before any data points are observed.
(i-k) show samples from the posterior, after the data has been observed. Note that the same data is used for all of the plots.}
\label{fig:composition_tutorial}
\end{figure}





Venture includes the primitive \texttt{make\_gp}, which takes as arguments a
unary function \texttt{mean} and a binary (symmetric, positive-semidefinite)
function \texttt{cov} and produces a function $\gtt$ distributed as a Gaussian
process with the supplied mean and covariance.  For example, a function $\gtt
\sim \GP(0,\,\SE)$, where $\SE$ is a squared-exponential covariance
\[ \SE(x, x') = \sigma^2 \exp\pn{\frac{(x-x')^2}{2\ell}} \]
with $\sigma=1$ and $\ell=1$, can be instantiated as follows:
\begin{lstlisting}[language=Venture]
assume zero = make_const_func( 0.0)
assume se = make_squaredexp( 1.0, 1.0)
assume g  = make_gp( zero, se)
\end{lstlisting}
There are two ways two view $\gtt$ as a ``random function.'' In the first view,
the \texttt{assume} directive that instantiates $\gtt$ does not use any
randomness---only the subsequent calls to $\gtt$ do---and coherence constraints
are upheld by the interpreter by keeping track of which evaluations of $\gtt$
exist in the current execution trace.  Namely, if the current trace contains evaluations
of $\gtt$ at the points $x_1,\ldots,x_N$ with return values $y_1,\ldots,y_N$,
then the next evaluation of $\gtt$ (say, jointly at the points $x_{N+1}, \ldots,
x_{N+n}$) will be distributed according to the joint conditional distribution
\[
  P\pn{\big.
    \texttt(\gtt\ x_{N+1}\texttt), \ldots, \texttt(\gtt\ x_{N+n}\texttt)
    \mvert
    \texttt(\gtt\ x_i\texttt) = f_i \text{ for $i=1,\ldots,N$}}.
\]
In the second view, $\gtt$ is a randomly chosen deterministic function, chosen
from the space of all deterministic real-valued functions; in this view, the
\texttt{assume} directive contains \emph{all} the randomness, and subsequent
invocations of $\gtt$ are deterministic.  The first view is procedural and is
faithful to the computation that occurs behind the scenes in Venture.  The
second view is declarative and is faithful to notations like ``$g \sim P(g)$''
which are often used in mathematical treatments.  Because a model program could
make arbitrarily many calls to $\gtt$, and the joint distribution on the return
values of the calls could have arbitrarily high entropy, it is not
computationally possible in finite time to choose the entire function $\gtt$ all
at once as in the second view.  Thus, it stands to reason that any
computationally implementable notion of ``nonparametric random functions'' must
involve incremental random choices in one way or another, and \ac{GP}s
in Venture are no exception.


\section{Venture GPs}

The Venture procedure \texttt{make-gp} takes as input a mean function and a covariance function, and outputs a procedure for sampling from a Gaussian process.
In effect, each call to this procedure samples from \eqref{eq:gpsampler} conditioned on the return values of all previous samples.
\texttt{make-gp} allows us to perform GP inference in Venture with only a few lines of code.
We can concisely express a wide variety of GPs: simple smoothing with fixed hyper-parameters, or a prior on hyper-parameters, or a custom covariance function.
Inference on hyper-parameters can be performed using Venture's built-in MH operator or a custom inference strategy.

Venture code to create and sample from a GP with a smoothing kernel and hyperparameters is shown in Listing \ref{alg:gpNeal}.
\begin{minipage}{\linewidth}
\small
\belowcaptionskip=-10pt
\begin{lstlisting}[frame=single,mathescape,label=alg:gpNeal,basicstyle=\selectfont\ttfamily]
// HYPER-PARAMETERS
assume log_sf = tag('hyper, log(gamma(1,1))))
assume log_l = tag('hyper, log(gamma(1,1))))

// COVARIANCE FUNCTION
assume se = make_squaredexp(log_sf, log_l)

// MAKE GAUSSIAN PROCESS
assume gp = make_gp( 0, se)

// INCORPORATE OBSERVATIONS
observe gp(array x[1],...x[n])= array(y[1],...,y[n])

// INFER HYPER-PARAMETERS
infer mh('hyper, one, 1)))

\end{lstlisting}
\end{minipage}


The first two lines declare the hyper-parameters.
We tag both of them to belong to the ``scope'' \texttt{'hyper}.
These tags are supplied to the inference program (in this case, MH) to specify on which random variables inference should be done.
In this \paperOrChapter, we use MH inference throughout.
Scopes may be further subdivided into blocks, on which block proposals can be made.
In this \paperOrChapter we do not use block proposals; MH inference is done on one variable at a time.

The ASSUME directives describe the GP model: \texttt{sf} and \texttt{l} (corresponding to $\sigma$ and $\ell$) are drawn from independent $\Gamma(1,3)$ distributions.
The squared exponential covariance function can be defined outside the Venture code in a conventional programming language (e.g. Python) and imported as a foreign SP.
In that way, the user can define custom covariance functions using his or her language and libraries of choice, without having to port existing code into Venture's modelling language.
In the above, the factory function \texttt{make-se}, which produces a squared exponential function with the supplied hyperparameters, is imported from Python (we have omitted the Python code).
In the next line \texttt{make-se} is used to produce a covariance function \texttt{SE}, whose (random) hyperparameters are \texttt{l} and \texttt{sf}.
Finally, we declare \texttt{GP} to be a Gaussian process with mean zero and covariance function \texttt{SE}.





{\color{red} I don't know what these two paragraphs mean --Ben}

In the case where hyper-parameters are unknown they can be found deterministically by optimizing the marginal likelihood using a gradient based optimizer. Non-deterministic, Bayesian representations of this case are also known~\citep{neal1997monte}. 

We have already implemented this in listing \ref{alg:gpNeal}. We draw the hyper-parameters from a $\Gamma$-prior for a Bayesian treatment of hyper-parameters. This is simple using the build in stochastic procedure that simulates drawing samples from a gamma distribution.
The program gives rise to a Bayesian representation of GPs, which we will explore in the following.




\subsection{Gaussian process memoization: \gpmem}
Memoization is the practice of storing previously computed values of a function so that future calls with the same inputs can be evaluated by lookup rather than re-computation.
To transfer this idea to probabilistic programming, we now introduce a language construct called a
\emph{statistical memoizer}.  Suppose we have a function $\ftt$ which can be evaluated 
but we wish to learn about the behavior of $\ftt$ using as
few evaluations as possible.  The statistical memoizer, which here we give the
name \gpmem, was motivated by this purpose.  It produces two outputs:
\[ \ftt \xrightarrow{\gpmem} (\ftt_{\text{compute}}, \ftt_\emu). \]
The function $\ftt_{compute}$ calls $\ftt$ and stores the output in a memo
table, just as traditional memoization does.  The function $\ftt_\emu$ is
an online statistical emulator which uses the memo table as its training
data.  A fully Bayesian emulator, modelling the true function $\ftt$ as a
random function $f \sim P(f)$, would satisfy
\[
\texttt{(}\ftt_\emu\ \xtt_1\ \ldots\ \xtt_k\texttt{)}
\sim
P\pn{
  f(\xtt_1), \ldots, f(\xtt_k)
  \mvert
  \text{$f(\xtt) = \texttt{(f x)}$ for each $\xtt$ in memo table}
}.
\]
Different implementations of the statistical memoizer can have
different prior distributions $P(f)$; in this paper, we deploy a \ac{GP} 
prior (implemented as \texttt{gpmem} below).  Note that we require the ability
to sample $\ftt_\emu$ jointly at multiple inputs because the values of
$f(\xtt_1),\ldots,f(\xtt_k)$ will in general be dependent.


\begin{figure}
\centering
\begin{tabular}{ll}\small
% line 1
\begin{lstlisting}[mathescape,escapechar=\#]
assume (f_compute f_emu) =  gpmem( f)
sample (f_emu( array( -20, $\cdots$, 20)) 

             #\raisebox{-0.2\height}{\includegraphics[height=.5cm]{figs/oneline.png}}# $\sim \mathcal{N}(\mu,\mathbf{K})$
\end{lstlisting}
& \raisebox{-0.5\height}{\includegraphics[height=2.5cm]{figs/slide1_0pred.png}} \\ \hline
% line 2
\begin{lstlisting}[mathescape,escapechar=\#]
predict f_compute( 12.6)

sample (f_emu( array( -20, $\cdots$, 20)) 

\end{lstlisting}
 &  \raisebox{-0.5\height}{\includegraphics[height=2.5cm]{figs/slide1_1pred.png}}  \\ \hline
% line 3
 \begin{lstlisting}[mathescape,escapechar=\#]
predict f_compute( -6.4)

sample (f_emu( array( -20, $\cdots$, 20)) 
  
\end{lstlisting}
 &   \raisebox{-0.5\height}{\includegraphics[height=2.5cm]{figs/slide1_2pred.png}}
\end{tabular}
%\put(-245,10){\line(1,0){200}}
\put(-33,77){\color{ForestGreen}\thicklines \vector(0,-1){15}}
\put(-96,6){\color{ForestGreen}\thicklines \vector(0,-1){15}}
\put(-48,-63){\thicklines \vector(0,-1){15}}
\put(-84,-43){\thicklines \vector(0,-1){15}}
\put(-73,-68){\thicklines \vector(0,1){15}}
\put(-460,110){(b)}
\put(-460,45){(c)}
\put(-460,-10){(d)}
\put(-460,-75){(e)}
\put(-460,-130){(f)}
\caption{\small \gpmem\ tutorial. The top shows a schematic of \gpmem.
  \texttt{f\_compute} probes an outside resource.
  This can be expensive (top left).
  Every probe is memoized and improves the emulator. Below the schematic we see the evolution
  of \gpmem's state of believe of the world given certain Venture
  directives. On the right, we depict the true function (blue), samples from the emulator (red) and incorporated observations (black).}
\label{fig:gpmem_tutorial}
\end{figure}

% Panel 1
We explain how \gpmem, the statistical memoizer with \ac{GP}-prior, works using a simple tutorial
(Fig. \ref{fig:gpmem_tutorial}). 
The top panel (Fig. \ref{fig:gpmem_tutorial}, (a)) of this figure sketches the schematic of \gpmem.
$\ftt$ is the external process that we memoize. It can be evaluated using resources that potentially come
from outside of Venture.  
We feed this function into \gpmem\ alongside
a parameterised kernel $\Ktheta$.  
In this example, we make the qualitative assumption of $f$ being smooth, and define
$\Ktheta$ to be a squared-exponential covariance function:
\[
\Ktheta = \text{SE} = \sigma^2 \exp(-\frac{(x-x^\prime)^2}{2\ell^2}).
\]
The hyper-parameters $\thetabf$ for this kernel are sampled from a 
prior distribution which is depicted in the top right box.
Note that we annotate $\bm{\theta}=\{\texttt{sf},\texttt{l}\}$ for subsequent
inference as belonging to (i) the scope "hyper-parameter" and (ii) blocks 0 and 1 respectively.

\gpmem\ implements a memoization table, where all previously
computed function evaluations are stored. We also initialize a \ac{GP}-prior that
will serve as our statistical emulator:
\[
P(f_{emu}(x) \mid \mathbf{x}_{past},\fbf)\sim \mathcal{N}(\mu(\mathbf{x}),\mathbf{K}_\theta\big(\mathbf{x},\mathbf{x})\big)
\]
where 
\[
P(f_{emu}(x) \mid \mathbf{x}_{past},\fbf) = \hat\fbf 
\]
under the traditional \ac{GP} perspective\footnote{Note that we write $\fbf_*$ instead of $\hat\fbf$ if $\xbf_{past}=\{\}$ and $\fbf=\{\}$.}.
All value pairs stored in the memoization table ($\text{memo table} = (\mathbf{x}_{past},\fbf)$) are incorporated as observations of
the \ac{GP}.
The emulator allows us to make probabilistic predictions on function evaluations, estimating
$P(\hat\fbf\mid \fbf)$ with $\mathcal{N}(\hat{\bm{\mu}},\hat\Kbf)$.
We simply feed the regression input
into the emulator and output a predictive posterior Gaussian distribution determined by the \ac{GP} and
the memoization table.

% Panel 2
We can either define the function f that serves as as input for \gpmem\
 natively in Venture
(as shown in the Fig. \ref{fig:gpmem_tutorial} (b)) or we interleave Venture with foreign code. 
This can be useful when $\ftt$ is computed with the help of outside resources.
Before making any observations or calls to $\ftt$
we can sample from the prior at the inputs from -20 to 20 using the emulator:
    \begin{lstlisting}
    assume (f_compute, f_emu) = gpmem(f, K))

    sample f_emu(array(-20, ..., 20))
    \end{lstlisting}
where the second line corresponds to:
\[ 
\fbf_* \sim \mathcal{N}\Bigg(0,K\bigg(
\bmat{
-20 \\
\cdots \\
20
},
\bmat{
-20 \\
\cdots \\
20
}
\bigg)
\Bigg).
\]
% Panel 3
In Fig. \ref{fig:gpmem_tutorial} (c), we probe the external function $\ftt$ at point 12.6 and memoize it's result by calling 
   \begin{lstlisting}
    predict f_compute (12.6).
    \end{lstlisting}
When we subsequently sample from the emulator, that is compute the $\hat\fbf$ at the input
$\hat\xbf= \bmat{-20, \cdots, 20}^\top$. We see how the posterior which shifts from uncertainty to near certainty close to the input 12.6.

% Panel 4
We can repeat the process at a different point (probing point -6.4 in Fig.
\ref{fig:gpmem_tutorial} (d)) to see that we gain certainty about another part of the curve. 

% Panel 5
We can add information to $\texttt{f}_\text{emu}$ about presumable value pairs of $\ftt$ without calling $\texttt{f}_\text{compute}$
(Fig. \ref{fig:gpmem_tutorial} (e)).
If a friend tells us the value of $\ftt$ we can call observe to store this information in the incorporated observations for $\texttt{f}_\text{emu}$ only:
    \begin{lstlisting}
    observe f_emu( -3.1) = 2.60.
    \end{lstlisting}
We have this value pair now available for the computation $\hat\fbf$. 
For sampling with the emulator, the effect is the same as calling predict with the $\texttt{f}_\text{compute}$.
However, we can imagine at least one scenario where such as distinction in the treatment of observations 
is beneficial. Let us say we do not only have the real function available but also a domain expert with knowledge 
about this function.
This expert could tell us what the value is at a given input.
Potentially, the value provided by the expert could disagree with the value computed with $\ftt$ for example 
due to different levels of observation noise. 

% Panel 6
Finally, we can update our posterior by inferring the posterior over hyper-parameter values $\thetabf$.
We can take 50 \ac{MH} steps by calling that where we take an \ac{MH} steps for either of the two blocks for the scope "hyper-parameter": 
   \begin{lstlisting}
    infer mh( 'hyper-parameter, one, 50).
    \end{lstlisting}
The newly inferred hyper-parameters allow us now adequately reflect uncertainty
about the curve given all incorporated observations (compare
Fig. \ref{fig:gpmem_tutorial}, bottom panel (f) on  the right with the samples
before inference, one panel above (e)).


\subsection{A Bayesian interpretation}
\begin{figure}
        \centering
        \begin{subfigure}[b]{0.49\textwidth} \centering
              \includegraphics[height=7.5cm]{figs/MAP_linear.png}
              \caption{MAP}
                \label{fig:maplin}
        \end{subfigure}%
        \begin{subfigure}[b]{0.49\textwidth} \centering
            \includegraphics[height=7.5cm]{figs/MH_linear.png}
            \caption{MH}
                \label{fig:mhlin}
        \end{subfigure}%
    
        \caption{Left, MAP estimate, right Metropolis Hastings}\label{fig:BayesVSFreq}
\end{figure}





\subsubsection{Data modelling as a special case of \gpmem}
\label{sec:special-case-gpmem}
We implement \gpmem\ by memoizing a target procedure in a wrapper that remembers previously computed values.
This comes with interesting implications:
from the standpoint of computation, a data set of the form $\{(x_i, y_i)\}$ can be thought of as a function $y = f_\restr(x)$, where $f_\restr$ is restricted to only allow evaluation at a specific set of inputs $x$ (Alg. \ref{alg:frest}).
\begin{algorithm}

\caption{Restricted Function}

\begin{algorithmic}[1]

\Function{$f_\restr$}{$x$}
\If{$x \in \mathcal{D}$}
 \State \Return $\mathcal{D}[x]$
\Else
 \State raise Exception('Illegal input')
\EndIf 
\EndFunction

\end{algorithmic}
  \label{alg:frest}


\end{algorithm}
Modelling the data set with a \ac{GP} then amounts to trying to learn a smooth function $f_\emu$ (``emu'' stands for ``emulator'') which extends $f$ to its full domain.
Indeed, if $f_\restr$ is a foreign procedure made available as a black-box to Venture, whose secret underlying pseudo code is in Alg. \ref{alg:frest},
then the \texttt{observe} code can be rewritten using \gpmem\ as in Listing \ref{alg:venfrest} (where here the data set \texttt{D} has keys \texttt{x[1]},\ldots,\texttt{x[n]}):
\begin{mdframed}
\begin{minipage}{0.9\linewidth}
\small
\begin{lstlisting}[mathescape,label=alg:venfrest,basicstyle=\selectfont\ttfamily,numbers=none,caption={Observation with \gpmem},escapechar=\#]
#\linenumber{1}#assume (f_compute f_emu) =  gpmem( f_restr)
#\linenumber{2}#for i=1 to n:
#\linenumber{3}#  predict f_compute( x[i])
#\linenumber{4}#  infer mh(quote(hyper-parameters), one, 100)
#\linenumber{5}#sample (f_emu( array( 1, 2, 3))
\end{lstlisting}
\end{minipage}
\end{mdframed}


\vspace{0.2cm}
This rewriting has at least two benefits: (i) readability (in some cases), and (ii) amenability to active learning.
As to (i), the statistical code of creating a Gaussian process is replaced with a memoization-like idiom, which will be more familiar to programmers.
As to (ii), when using \gpmem, it is quite easy to decide incrementally which data point to sample next: for example, the loop from \texttt{x[1]} to \texttt{x[n]} could be replaced by a loop in which the next index \texttt{i} is chosen by a supplied decision rule.
In this way, we could use \gpmem\ to perform online learning using only a subset of the available data.

\begin{figure}
\centering
\begin{tabular}{ll}\small
% line 1
\begin{lstlisting}[mathescape,escapechar=\#]
assume (f_compute f_emu) =  gpmem( f)
sample (f_emu( array( -20, $\cdots$, 20)) 

             #\raisebox{-0.2\height}{\includegraphics[height=.5cm]{figs/oneline.png}}# $\sim \mathcal{N}(\mu,\mathbf{K})$
\end{lstlisting}
& \raisebox{-0.5\height}{\includegraphics[height=2.5cm]{figs/slide1_0pred.png}} \\ \hline
% line 2
\begin{lstlisting}[mathescape,escapechar=\#]
predict f_compute( 12.6)

sample (f_emu( array( -20, $\cdots$, 20)) 

\end{lstlisting}
 &  \raisebox{-0.5\height}{\includegraphics[height=2.5cm]{figs/slide1_1pred.png}}  \\ \hline
% line 3
 \begin{lstlisting}[mathescape,escapechar=\#]
predict f_compute( -6.4)

sample (f_emu( array( -20, $\cdots$, 20)) 
  
\end{lstlisting}
 &   \raisebox{-0.5\height}{\includegraphics[height=2.5cm]{figs/slide1_2pred.png}}
\end{tabular} 
%\put(-245,10){\line(1,0){200}}
\put(-39,30){\thicklines \vector(0,-1){20}}
\put(-113,-42){\thicklines \vector(0,-1){20}}
\caption{\gpmem\ tutorial}
\end{figure}

\subsubsection{The efficacy of learning hyperparameters}
We can apply \gpmem\ for regression in a hierarchical Bayesian setting
(Fig. \ref{fig:neal_tutorial}).  
%%%%%%%%%%%%%%%%%%%%%%%%%%%%%%%%%%%%%%%%%%%%%%%%%%%%%%%%%%%%
\begin{figure}
\renewcommand{\arraystretch}{0.1}% Tighter
\centering \footnotesize
\begin{tikzpicture}
\node[] (start) {};
\node[left=1cm of start] (sfsigma) {\includegraphics[height=2cm]{figs/hypers_sf_sigma.png}};
\node[right=1cm of start] (sfell) {\includegraphics[height=2cm]{figs/hypers_sf_ell.png}};
\node[above=1cm of start] (hyper) {(a) $P(\bm{\theta} = \{\ell,sf,\sigma\} \mid
\mathbf{D},\Krv)$};
\node[left=0.cm of sfsigma] (ell) {$\ell$};
\node[below=0.cm of sfsigma] (sf1) {$sf$};
\node[left=0.cm of sfell] (sigma) {$\sigma$};
\node[below=0.cm of sfell] (sf2) {$sf$};






%%%%%%%%%%%%%%%%%%%%%%%%%%%%%%%%%%%%%%%%%%%%%%%%%%%%%%%%%%%%
%%%%%%%%%%%%  Code             %%%%%%%%%%%%%%%%%%%%%%%%%%%%%
%%%%%%%%%%%%%%%%%%%%%%%%%%%%%%%%%%%%%%%%%%%%%%%%%%%%%%%%%%%%

\node[below=1.0cm of start] (b_code){
\footnotesize\begin{lstlisting}[mathescape,escapechar=\@]
// Define data and look-up function
define data      = array(array(-1.87, 0.13),@\ldots@, array(1.67,0.81));
assume f_look_up = proc(index) {lookup(data, index)};
\end{lstlisting}
};
\node[below=-0.65cm of b_code,xshift=-0.58cm] (c_code){
\footnotesize\begin{lstlisting}[mathescape,escapechar=\@
]
// Initialize hyper-priors
assume alpha_sf @$\sim$@ gamma(5,1)                 #"hyperhyper";
assume beta_sf  @$\sim$@ gamma(5,1)                 #"hyperhyper";
assume alpha_l  @$\sim$@ gamma(5,1)                 #"hyperhyper";
assume beta_l   @$\sim$@ gamma(5,1)                 #"hyperhyper";@\vspace{1mm}@
assume sf       @$\sim$@ gamma(alpha_sf, beta_sf))) #"hyper";
assume l        @$\sim$@ gamma(alpha_l, beta_l)))   #"hyper";
assume sigma    @$\sim$@ gamma(5,1))                #"hyper"; 
\end{lstlisting}
};
% (d) & (e)
\node[below=-0.65cm of c_code,xshift=-0.65cm] (d_code){
\footnotesize\begin{lstlisting}[mathescape,escapechar=\@]
// Initialize covariance function
assume se = @\se@;
assume wn = @\wn@;
assume composite_covariance = gp_cov_sum(se, wn);
// Create a prober and emulator using gpmem
assume (f_compute, f_emu) =
    gpmem(f_look_up, composite_covariance);@\vspace{1mm}@
sample f_emu(array(-2, $\cdots$, 2));
\end{lstlisting}
};
% (f)
\node[below=0.16cm of d_code,xshift=-0.18cm] (e_code){
\footnotesize\begin{lstlisting}[mathescape,escapechar=\@]

// Observe all data points
for (x,y) in data {
                   observe f_emu(x) = y};
// Or: probe all data points
for (x,_) in data {
                   predict f_compute(x)};@\vspace{1mm}@
sample f_emu(array(-2, $\cdots$, 2));
\end{lstlisting}
};
% (g)
\node[below=-0.0cm of e_code,xshift=-0.77cm] (f_code){
\footnotesize\begin{lstlisting}[mathescape,escapechar=\@]
// Metropolis-Hastings
infer repeat(100, do(
                mh(#"hyperhyper", steps=1),
                mh(#"hyper",      steps=3)));@\vspace{1mm}@
sample f_emu(array(-2, $\cdots$, 2));
\end{lstlisting}
};

\node[below=0.5cm of f_code,xshift=0.3cm] (g_code){
\footnotesize\begin{lstlisting}[mathescape,escapechar=\@]
// Optimization
infer gradient-ascent(#"hyper", steps=10);

sample f_emu(array(-2, $\cdots$, 2));
\end{lstlisting}
};



% samples/curve images
\node[below=1.1cm of c_code,xshift=6cm] (e_pic){
\includegraphics[height=3.4cm]{figs/neal_example_before_observation.jpg}
};
\node[below=-0.4cm of e_pic] (f_pic){
\includegraphics[height=3.4cm]{figs/neal_example_after_observation.jpg}
};
\node[below=-0.4cm of f_pic] (g_pic){
\includegraphics[height=3.4cm]{figs/neal_Bayesian.jpg}
};
\node[below=-0.4cm of g_pic] (h_pic){
\includegraphics[height=3.4cm]{figs/neal_example_map_inference_alpha0p01_iter15.jpg}
};

% horizontal lines
% (a)
\draw 
  ([xshift=-0.2cm,yshift=-0.6cm]b_code.north west) --
([xshift=0.8cm,yshift=-0.6cm]b_code.north east); 
% (b)
\draw 
  ([xshift=-0.2cm,yshift=0.08cm]b_code.south west) -- ([xshift=0.8cm,yshift=0.08cm]b_code.south east);   
% (c)
 \draw 
  ([xshift=-0.2cm,yshift=-3.17cm]b_code.south west) --
([xshift=0.8cm,yshift=-3.17cm]b_code.south east); 
% (d)
 \draw 
  ([xshift=-0.2cm,yshift=-4.78cm]b_code.south west) --
([xshift=0.8cm,yshift=-4.78cm]b_code.south east); 
% (e)
 \draw 
  ([xshift=-0.2cm,yshift=-7.7cm]b_code.south west) -- ([xshift=0.8cm,yshift=-7.7cm]b_code.south east); 
% (f)
 \draw 
  ([xshift=-0.2cm,yshift=-10.8cm]b_code.south west) -- ([xshift=0.8cm,yshift=-10.8cm]b_code.south east); 
% (g)
 \draw 
  ([xshift=-0.2cm,yshift=-14.1cm]b_code.south west) -- ([xshift=0.8cm,yshift=-14.1cm]b_code.south east); 

% labels
\node[below=1.9cm of start,xshift=-6.5cm] (b) {(b)};
\node[below=1.4cm of b] (c) {(c)};
\node[below=1.9cm of c] (d) {(d)};
\node[below=1.7cm of d] (e) {(e)};
\node[below=2.6cm of e] (f) {(f)};
\node[below=2.9cm of f] (g) {(g)};
\node[below=2.6cm of g] (h) {(h)};
\end{tikzpicture}


\put(-410,230){(b)}
\put(-410,170){(c)}
\put(-410,90){(d)}
\put(-410,-10){(e)}
\put(-410,-110){(f)}
\put(-410,-210){(g)}
\caption{Regression with outliers and hierarchical prior structure}
\label{fig:neal_tutorial}
\end{figure}
%%%%%%%%%%%%%%%%%%%%%%%%%%%%%%%%%%%%%%%%%%%%%%%%%%%%%%%%%%%%
In a Bayesian treatment of  hyper-parameter learning for \ac{GP}s,
we can write the posterior probability of the hyper-parameters of a GP  (Fig.
\ref{fig:neal_tutorial}, (a)) given covariance function $\mathbf{K}$ as:
%%%%%%%%%%%%%%%%%%%%%%%%%%%%%%%%%%%%%%%%%%%%%%%%%%%%%%%%%%%%
\begin{equation}
\label{eq:hyperProbability}
P(\bm{\theta}=\{sf,\ell,\sigma\} \mid \mathbf{D,K}) = \frac{P(\mathbf{D} \mid \bm{\theta}, \mathbf{K})P(\bm{\theta} \mid  \mathbf{K})}{P(\mathbf{D} \mid \mathbf{K})}
\end{equation}
%%%%%%%%%%%%%%%%%%%%%%%%%%%%%%%%%%%%%%%%%%%%%%%%%%%%%%%%%%%%
where $\mathbf{D} = \{\xbf, \fbf\}$ is a training data set. Since we can apply
\gpmem\ to any process or procedure, it can be used in situations where only a data
set is available via a look-up function $\texttt{f\_look\_up}$.
In fact, we demonstrate \gpmem's application to regression using an example where
the data was generated by a function which is not available (Fig.
\ref{fig:neal_tutorial} (b)).
This function, $f_\text{true}$, is taken from a paper on the
treatment of outliers with hierarchical Bayesian hyper-priors for
\ac{GP}s~\citep{neal1997monte}:
%%%%%%%%%%%%%%%%%%%%%%%%%%%%%%%%%%%%%%%%%%%%%%%%%%%%%%%%%%%%
\begin{equation}
f_\text{true}(x) =  0.3 + 0.4 x + 0.5 \sin(2.7x) + \frac{1.1}{(1+ x^2)} + \eta
\;\;\; with\;\;\eta \sim \mathcal{N}(0,\sigma).
\end{equation}
%%%%%%%%%%%%%%%%%%%%%%%%%%%%%%%%%%%%%%%%%%%%%%%%%%%%%%%%%%%%
We synthetically generate outliers by setting $\sigma = 0.1$ in $95\%$ of the cases and to $\sigma = 1$ in the remaining cases. 
Instead of accessing the $f_\text{true}$ directly, we are accessing the $\texttt{data}$ in form of
a a two dimensional $\texttt{array}$ with $\texttt{f\_look\_up}$.

We parameterize $\Kbf_{\bm{\theta}}$ with $\bm{\theta}=\{sf,\ell,\sigma\}$.
For these hyper-parameters, Neals work suggests a hierarchical system for
for hyper-parameterization.
Here, we draw hyper-parameters from $\Gamma$ distributions:
%%%%%%%%%%%%%%%%%%%%%%%%%%%%%%%%%%%%%%%%%%%%%%%%%%%%%%%%%%%%
\begin{equation}
\label{eq:hyper-ell}
\ell^{(t)} \sim \Gamma(\alpha_1,\beta_1),\;\sigma^{(t)} \sim \Gamma(\alpha_2,\beta_2)
\end{equation} 
%%%%%%%%%%%%%%%%%%%%%%%%%%%%%%%%%%%%%%%%%%%%%%%%%%%%%%%%%%%%
and in turn sample the $\alpha$ and $\beta$ from $\Gamma$ distributions as well:
%%%%%%%%%%%%%%%%%%%%%%%%%%%%%%%%%%%%%%%%%%%%%%%%%%%%%%%%%%%%
\begin{equation}
\label{eq:hyper-alpha}
\alpha_1^{(t)} \sim \Gamma(\alpha^1_{\alpha},\beta^1_{ \alpha } ),\; \alpha_2^{(t)} \sim \Gamma(\alpha^2_{\alpha},\beta^2_{\alpha}),\cdots
\end{equation}
%%%%%%%%%%%%%%%%%%%%%%%%%%%%%%%%%%%%%%%%%%%%%%%%%%%%%%%%%%%%
We model this in Venture as illustrated in Fig. \ref{fig:neal_tutorial} (c),
using the build-in \ac{SP} $\texttt{gamma}$. Note that we omit the prior distributions of (\ref{eq:hyper-alpha}) due limited space in the tutorial.

In Fig. \ref{fig:neal_tutorial}, panel (d), we see that $\mathbf{K}_{\bm{\theta}}$
is defined as a composite covariance function. It is the sum ($\texttt{add\_funcs}$) of
an SE kernel ($\texttt{make\_squaredexp}$) and a white noise (WN, Appendix A)
kernel which is implemented with $\texttt{make\_whitenoise}$\footnote{Note
that in Neal's work \citeyearpar{neal1997monte} the sum of an SE
plus a constant kernel is used. We use a WN kernel for illustrative purposes
instead.}. 
We then initialize \gpmem\ feeding it with $\texttt{composite\_covariance}$ and the data
look-up function $\texttt{f\_look\_up}$. 
We sample from the prior $\fbf_*$ with random parameters $\texttt{sf,l}$ and $\texttt{sigma}$ and 
without any observations available.
We depict those samples on the right (red), alongside the true function that generated the data (blue) and
the data points we have available in the data set (black).

We can incorporate observations using both \texttt{observe} and \texttt{predict} (Fig. \ref{fig:neal_tutorial} (e)).
When we subsequently sample $\hat\fbf$ from the emulator with
$\mathcal{N}(\hat{\bm{\mu}},\hat\Kbf)$, we can see that the \ac{GP} posterior incorporates knowledge 
about the $\texttt{data}$. Yet, the hyper-parameters $\texttt{sf,l}$ and $\texttt{sigma}$ are still
random, so the emulator does not capture the true underlying dynamics
($f_\text{true}$) of the \texttt{data} correctly. 

Next, we demonstrate how we can capture these underlying dynamics within only
100  nested \ac{MH} steps on the hyper-parameters to get a good approximation
for their posterior $\hat\fbf$ (Fig. \ref{fig:neal_tutorial} (f)).
We say nested because we first take two sweeps in the scope
$\texttt{hyperhyper}$ which characterizes (\ref{eq:hyper-alpha}) and then one
sweep on the scope $\texttt{hyper}$ which characterizes (\ref{eq:hyper-ell}).
This is repeated 100 using $\texttt{repeat( 100, do(}\cdots\;$.
Note that Neal devises an additional noise model and performs a large number of Hybrid-Monte Carlo and Gibbs steps to achieve this, whereas inference in Venture with \gpmem\ is merely one line of code. 

Finally, we can change our inference strategy altogether. If we decide that instead of
following a Bayesian sampling approach, we would like to perform empirical optimization,
we do this by only changing one line of code, deploying $\texttt{map}$ instead
of $\texttt{mh}$ (Fig. \ref{fig:neal_tutorial} (g)). 


\subsubsection{Broader applicability of \gpmem}\label{sec:gpmem-broader}
\input{sections/applicability.tex}
\subsection{Structure Learning}\label{sec:structurelearning}
Inductive learning of symbolic expressions for continuous-valued time series
data is a hard task which has recently been tackled using a greedy search over 
the approximate posterior of the possible kernel compositions for
\ac{GP}s~\citep{duvenaud2013structure,lloyd2014automatic}\footnote{\url{http://www.automaticstatistician.com/}}.

With \gpmem\ we can provide a fully Bayesian treatment of this, previously unavaible,
using a stochastic grammar  (see Fig. \ref{fig:schema}).
%%%%%%%%%%%%%%%%%%%%%%%%%%%%%%%%%%%%%%%%%%%%%%%%%%%%%%%%%%%%
\begin{figure}
\centering
\usetikzlibrary{arrows, decorations.markings}
\usetikzlibrary{trees}

\tikzstyle{level 1}=[level distance=1cm, sibling distance=1.5cm]
\tikzstyle{level 2}=[level distance=1cm, sibling distance=1.1cm]

% Define styles for operators and leafs
\tikzstyle{operator} = [draw=none,circle, minimum width=1pt]
\tikzstyle{end} = [circle, minimum width=3pt,fill, inner sep=0pt]
% for double arrows a la chef
% adapt line thickness and line width, if needed

\begin{subfigure}[b]{0.49\textwidth}\centering
\begin{tikzpicture}[thick]

 \node[] (base_kernels) {\small$\text{BK}=\{k_{\bm{\theta}^1}^1,\cdots,k_{\bm{\theta}^m}^m\}$};
 \node[above=0.7cm of base_kernels] (theta) {\small$\bm{\theta}^* \sim P(\bm{\theta}^*)$};
 \node[draw,rectangle, below=1cm of start, text width =6.0cm, text 
height=4.1cm,align=center] (grammar)
{};


 \node[draw,dashed,rectangle, below=-3.8cm of grammar, text width
=5.3cm,align=center] (subset) {
$\;$Select Primitive Kernels\\% I can't believe that this is the simpliest way to
% get spacing right 
{\raggedright
\footnotesize$\Sbf  \sim P(\Sbf =
\{k_{\bm{\theta}^i}^i,\cdots,k_{\bm{\theta}^n}^n\} \mid \text{BK})$
}
};
 \node[draw,rectangle,dashed, below=0.7cm of subset, text width
=4.35cm,align=center]
(composition_procedure) {
$\;$Kernel Composer\\ % I can't believe that this is the simpliest way to
% get spacing right 
{\raggedright
\footnotesize$\;\;\;\bm{\Omega} \sim P(\bm{\Omega} \mid \Sbf)$ \\

\footnotesize$k_{\bm{\theta}} \sim P(\Krv \mid \bm{\Omega},\Sbf),\;\,\bm{\theta}\subseteq\bm{\theta}^*$
}
};

\node[above=0.0cm of subset]{\centering \bf Stochastic Grammar}; 

 \node[draw,rectangle,below=1cm of grammar] (gpmem) {\texttt{gpmem}};

 \node[draw,rectangle,below=1cm of gpmem] (f) {Data Generation};
\node[left of=f,xshift=-1.5cm] (x) {$\xbf$};
\node[right of=f,xshift=1.5cm] (y) {$\ybf$};

 \node[below=0.1cm  of grammar,xshift=0.5cm] (k)
{\small $k_{\thetabf}$};


\node[below=0.1cm  of gpmem,xshift=0.6cm] (gp) {
\small$f_{emu}$};

% 1st pass: draw arrows
  \draw[thick,->] (base_kernels) -- (grammar);
 
  \draw[thick,->] (theta) -- (base_kernels);
  \draw[thick,->] (grammar) -- (gpmem);
  \draw[thick,->] (gpmem) -- (f);
 \draw[thick,->] (x) -- (f);
 \draw[thick,->] (f) -- (y);
 \draw[thick,dashed,->] (subset) -- node[right]{\footnotesize $\;\Sbf \subseteq
\text{BK}$} (composition_procedure);
  % Note: If you have no branches, the 2nd pass is not needed
\end{tikzpicture}\vspace{2mm}
\caption{} 
\end{subfigure}
\addtocounter{subfigure}{2}
\begin{subfigure}[b]{0.49\textwidth}\centering
\begin{tikzpicture}[grow=right, sloped]
\node[operator,yshift=5.2cm] {\small $+$}
    child {
        node[operator] {\small $+$}        
            child {
               node[operator] {\small $\times$}        
        child {
                node[operator, label=right:
                    {$\cdots$}] {}
                edge from parent
                node[above] {}
                node[below]  {}
            }
            child {
                node[end, label=right:
                    {LIN$_{\theta^4}$}] {}
                edge from parent
                node[above] {}
                node[below]  {}
            }
        edge from parent         
            node[above] {}
            node[below]  {}
            }
            child {
                node[end, label=right:
                    {WN$_{\theta^3}$}] {}
                edge from parent
                node[above] {}
                node[below]  {}
            }
            edge from parent 
            node[above] {}
            node[below]  {}
    }
    child {
        node[operator] {\small $\times$}        
        child {
                node[end, label=right:
                    {SE$_{\theta^2}$}] {}
                edge from parent
                node[above] {}
                node[below]  {}
            }
            child {
                node[end, label=right:
                    {SE$_{\theta^1}$}] {}
                edge from parent
                node[above] {}
                node[below]  {}
            }
        edge from parent         
            node[above] {}
            node[below]  {}
    };


\node[xshift=-1.4cm,yshift=5.2cm] (K) {$\Parse(\ktheta)=$}; 
\node[draw,rectangle,font=\tiny,text width = 6.7 cm, yshift=0.8cm] (simplification) {
SE $\times$ SE $\;\;\;\;\;\;\;\;\;\;\;\;\;\;\;\;\;\;\;\;\;\;\;\;\;\;\;\;\;\;\,\rightarrow$ SE\\
$\{$SE, PER, C, WN$\} \times$ WN $\;\;\;\;\,\rightarrow$ WN\\
$\{$SE, PER, C, WN, LIN$\} \times$ C $\rightarrow$ $\{$SE, PER, C, WN, LIN$\}$\\
LIN $+$ LIN $\;\;\;\;\;\;\;\;\;\;\;\;\;\;\;\;\;\;\;\;\;\;\;\;\;\;\;\rightarrow$ LIN\\
PER $+$ LIN $ \;\;\;\;\;\;\;\;\;\;\;\;\;\;\;\;\;\;\;\;\;\;\;\;\;\,\rightarrow$ LIN $+$ PER, $\cdots$\\
PER $\times$ LIN $\;\;\;\;\;\;\;\;\;\;\;\;\;\;\;\;\;\;\;\;\;\;\;\;\;\,\rightarrow$ LIN $\times$ PER,$\cdots$\\
};
\node[above=0.0cm of simplification](simplification_with){\small Simplificaton 
 with $\Simplify()$};

\node[yshift=-3cm] (Keq)
{$\text{Struct}(k)=\text{SE} +\text{WN} + \text{LIN} \times ( \cdots )$};
\node[yshift=-0.3cm](helper){};
\node[yshift=-2.7cm](helper2){};
\node[above=-0.1cm of
Keq](interpretation){\small$\kper \rightarrow \text{PER},\;\klin \rightarrow
\text{LIN},\;\kse \rightarrow \text{SE},\cdots$};
\node[above=0.0cm of interpretation](interpretation_with){\small Interpretation 
 with Struct()};
\node[draw, rectangle, below=0.0cm of interpretation_with,minimum width = 8cm, minimum height
=1.2cm] (structbox){};
\node[below=0.3cm of simplification,xshift=-0.2cm]{(c)};
\node[above=0.6cm of simplification_with,xshift=-0.2cm]{(b)};
\end{tikzpicture}
\caption{} 
\end{subfigure}

\begin{subfigure}[b]{0.99\textwidth}\centering
\begin{tabular}{cccc}
\multicolumn{4}{c}{{\bf Base Kernels} (BK)} \rule{0pt}{3ex} \\ 
\small LIN: Linearity &\small PER: Periodicity &\small SE: Smoothness &\small WN: White Noise \rule{0pt}{2ex} \\
\includegraphics[height=2cm]{figs/kernel/kernelLIN.png} & \includegraphics[height=2cm]{figs/kernel/kernelPER.png} & \includegraphics[height=2cm]{figs/kernel/kernelSE.png} & \includegraphics[height=2cm]{figs/kernel/kernelWN.png}\\
\end{tabular}
\begin{tabular}{cccc}
\multicolumn{4}{c}{\bf Composite Structure} \rule{0pt}{0ex}  \\ 
\small LIN + PER: &\small LIN $\times$ PER: &\small SE $\times$ PER: &\small LIN $\times$ LIN: \rule{0pt}{2ex} \\
\small Periodicity with Trend &\small Growing Amplitude &\small Local Periodicity&\small Quadratic \rule{0pt}{2ex} \\
\includegraphics[height=2cm]{figs/kernel/kernelLINplusPER.png} & \includegraphics[height=2cm]{figs/kernel/kernelLINtimesPER.png} & \includegraphics[height=2cm]{figs/kernel/kernelSEplusPER.png}& \includegraphics[height=2cm]{figs/kernel/kernelLINtimesLIN.png}\\
\end{tabular}
\caption{}
\end{subfigure}

\caption{\small(a) Bayesian GP structure learning. A set of
base kernels (BK) with priors on there hyper-parameters serve as hypothesis space
and is supplied as input to the stochastic grammar. The stochastic grammar has
two parts: (i) a sampler that selects a random set $\Sbf$  of primitive kernels from BK
and (ii) a kernel composer that combines the individual base kernels and generates
a composite kernel function
$k_{\thetabf}$. This serves as input for
\gpmem.  We observe value pairs $\xbf,\ybf$ of unstructured time series data on
the bottom of the schematic. (b) An example of composite kernel structure. We
use the Struct($k_{\thetabf}$) to compute a symbolic interpretation for
$k_{\thetabf}$. We use addition and multiplication to combine base kernels.
Base kernels and compositional kernels are shown in (c) alongside ther
interpretation with Struct().}\label{fig:schema}
\end{figure}
%%%%%%%%%%%%%%%%%%%%%%%%%%%%%%%%%%%%%%%%%%%%%%%%%%%%%%%%%%%%
This allows us to read an unstructured time series and automatically output a high-level,
qualitative description of it. The stochastic grammar takes a set of primitive base kernels 
    $\text{BK}=\{k_{\bm{\theta}^1}^1,\cdots,k_{\bm{\theta}^m}^m\}$
    of size $m$, including their corresponding parametrization
    $\thetabf^*=\{\thetabf^1,\cdots,\thetabf^m\}$ (Fig. \ref{fig:schema} (a) and (b))
We depict the input for the
stochastic grammar in Listing \ref{alg:base_kernels}.
\begin{mdframed}
\begin{minipage}{\linewidth}
\small
\belowcaptionskip=-10pt
\begin{lstlisting}[mathescape,label=alg:base_kernels,basicstyle=\selectfont\ttfamily,numbers=none,caption={Initialize
Base Kernels BK and $P(n)$},escapechar=\#]
#\linenumber{1}# // Initialize hyper-parameters
#\linenumber{2}#assume theta_1 = tag(scope="hyper-parameters", block=1, gamma(5,1));
#\linenumber{3}#assume theta_2 = tag(scope="hyper-parameters", block=2, gamma(5,1));
#\linenumber{4}#assume theta_3 = tag(scope="hyper-parameters", block=3, gamma(5,1));
#\linenumber{5}#assume theta_4 = tag(scope="hyper-parameters", block=4, gamma(5,1));
#\linenumber{6}#assume theta_5 = tag(scope="hyper-parameters", block=5, gamma(5,1));
#\linenumber{7}#assume theta_6 = tag(scope="hyper-parameters", block=6, gamma(5,1));
#\linenumber{8}#assume theta_7 = tag(scope="hyper-parameters", block=7, gamma(5,1));
#\linenumber{9}#
#\linenumber{11}# // Make kernels
#\linenumber{12}#assume lin = apply_function(make_linear, theta_1);
#\linenumber{13}#assume per = apply_function(make_periodic, theta_2, theta_3, theta_4);
#\linenumber{14}#assume se  = apply_function(make_squaredexp, theta_5, theta_6);
#\linenumber{15}#assume wn  = apply_function(make_noise, theta_7);
#\linenumber{16}#
#\linenumber{17}#// Initialize the set of primitive base kernels BK 
#\linenumber{18}#assume BK = list(lin, per, se, wn);
\end{lstlisting}
\end{minipage}
\end{mdframed}


We sample a random subset S of
the set of supplied base kernels. S is of size $n \leq m$. We write
%%%%%%%%%%%%%%%%%%%%%%%%%%%%%%%%%%%%%%%%%%%%%%%%%%%%%%%%%%%%
\[
\Sbf = \{K_{\bm{\theta}^i}^i,\cdots,K_{\bm{\theta}^n}^n\}
\sim P(\Sbf = \{K_{\bm{\theta}^i}^i,\cdots,K_{\bm{\theta}^n}^n\} \mid
\text{BK}) 
\]
%%%%%%%%%%%%%%%%%%%%%%%%%%%%%%%%%%%%%%%%%%%%%%%%%%%%%%%%%%%%
with
%%%%%%%%%%%%%%%%%%%%%%%%%%%%%%%%%%%%%%%%%%%%%%%%%%%%%%%%%%%%
\[
P(\Sbf = \{K_{\bm{\theta}^i}^i,\cdots,K_{\bm{\theta}^n}^n\}\mid \text{BK}) =
\frac{n!}{ \mid \Sbf = \{K_{\bm{\theta}^i}^i,\cdots,K_{\bm{\theta}^n}^n\}\mid !}.
\]
%%%%%%%%%%%%%%%%%%%%%%%%%%%%%%%%%%%%%%%%%%%%%%%%%%%%%%%%%%%%

BK is assumed to be fixed as the most general margin of our hypothesis space.
In the following, we will drop it in the notation.
The only building block that we are now missing is how to combine the sampled
base kernels into a compositional covariance function (see Fig. \ref{fig:schema}
(b)). For each interaction $i$, we
have to infer whether the data supports a local interaction or a global interaction,
chosing between one out of two algebraic operators
$\bm{\Omega}_i=\{+,\times\}$. The probability for all such decisions is given by a binomial distribution: 
\begin{equation}
P(\bm{\Omega} \mid \Sbf)= {n \choose r}  p_{+\times}^r (1 - p_{+\times})^{n-r}.
\end{equation}
We can write the marginal probability of a kernel function as 
%%%%%%%%%%%%%%%%%%%%%%%%%%%%%%%%%%%%%%%%%%%%%%%%%%%%%%%%%%%%
\begin{equation}
P(\Krv \mid \xbf,\ybf,\thetabf ) = \iint \limits_{\bm{\Omega},\Sbf}
P(\Krv \mid \xbf,\ybf,\thetabf,\bm{\Omega},\Sbf) \times P(\bm{\Omega} \mid \Sbf)\times
P(\Sbf)\; \text{\bf d}\bm{\Omega}\, \text{\bf d}\Sbf\,
\end{equation}
%%%%%%%%%%%%%%%%%%%%%%%%%%%%%%%%%%%%%%%%%%%%%%%%%%%%%%%%%%%%
with $\bm{\theta}\subseteq \bm{\theta}^*$ as implied by $S$.
For structure learning with \ac{GP} kernels, a composite kernel is
sampled from $P(\Krv)$ and fed into \gpmem. 
The emulator generated by \gpmem\ observes unstructured time series data.
Venture code for the probabilistic grammar is shown in Listing
\ref{alg:grammar}, code for inference with \gpmem\ in Listing
\ref{alg:structureVent}. 


\begin{mdframed}
\begin{minipage}{\linewidth}
\small
\belowcaptionskip=-10pt
\begin{lstlisting}[mathescape,label=alg:grammar,basicstyle=\selectfont\ttfamily,numbers=none,caption={
Stochastic Grammar},escapechar=\#]
#\linenumber{1}#// Select a random subset of the set possible primitive kernels (BK)
#\linenumber{2}#assume primitive_kernel_selection = tag("grammar", 0,
#\linenumber{3}#				 select_primitive_kernels(BK));
#\linenumber{4}#// Construct kernel composition with a composer procedure
#\linenumber{5}#assume kernel_composer = proc(l) {
#\linenumber{6}#  if (size(l) <= 1) {
#\linenumber{7}#    first(l)
#\linenumber{8}#  } else {
#\linenumber{9}#       if (bernoulli()) {
#\linenumber{10}#            add_funcs(first(l),  kernel_composer(rest(l)))
#\linenumber{11}#       } else {
#\linenumber{12}#            mult_funcs(first(l), kernel_composer(rest(l)))
#\linenumber{13}#    }
#\linenumber{14}#  }
#\linenumber{15}#};
#\linenumber{16}#
#\linenumber{17}#assume K = tag("grammar", 1,
	        	kernel_composer(primitive_kernel_selection));
\end{lstlisting}

\end{minipage}
\end{mdframed}

\begin{mdframed}
\begin{minipage}{\linewidth}
\small
\belowcaptionskip=-10pt
\begin{lstlisting}[mathescape,label=alg:structureVent,basicstyle=\selectfont\ttfamily,numbers=none,caption={\gpmem\
inference for structure
learning: },escapechar=\#]
#\linenumber{1}#// Apply gpmem 
#\linenumber{2}#assume (f_compute, f_emu) = gpmem(f_look_up, K);
#\linenumber{3}#// Probe all data points
#\linenumber{4}#for (n = 0; n < size(data); n++) { 
#\linenumber{5}#	predict f_compute(first(lookup(data, n)))};
#\linenumber{6}#// Perform inference
#\linenumber{7}#infer repeat(200, do(
#\linenumber{8}#	mh(scope="grammar", steps=1),
#\linenumber{9}#	mh(scope="hyper-parameters", steps=2)));
\end{lstlisting}

\end{minipage}
\end{mdframed}






Many equivalent covariance structures can be sampled due to covariance function algebra
and equivalent representations with different parameterization~\citep{lloyd2014automatic}.
To inspect the posterior of these equivalent structures we convert each kernel expression
into a sum of products and subsequently simplify. All base kernels can be found in Appendix A,
rules for this simplification can be found in Appendix B.
Kernel structure comes with symbolic interpretations. To read kernel function $k$
and apply the simplifications described above, we apply, Struct$(k)$. For
examle, we write
\[
\Struct(\klin)=\text{LIN},
\]
which translates a function into a symbolic expression, see Appendix C. 
%%%%%%%%%%%%%%%%%%%%%%%%%%%%%%%%%%%%%%%%%%%%%%%%%%%%%%%%%%%%%%%%%%%%%%%%%%%%%%%%%
%%%%%%%%%%%%%%%%%%%%%%%%%     Mauna result      %%%%%%%%%%%%%%%%%%%%%%%%%%%%%%%%%%
%%%%%%%%%%%%%%%%%%%%%%%%%%%%%%%%%%%%%%%%%%%%%%%%%%%%%%%%%%%%%%%%%%%%%%%%%%%%%%%%%
We defined a simple space of covariance structures in a way that allows us to produce results coherent with 
work presented in Automatic Statistician. The results are illustrated with two data sets.

\myparagraph{Mauna Loa  CO$_2$ data}
%%%%%%%%%%%%%%%%%%%%%%%%%%%%%%%%%%%%%%%%%%%%%%%%%%%%%%%%%%%%%%%%%%%%%%%%%%%%%%%%%
\begin{figure}
\centering
 \addtolength\abovedisplayskip{-1\baselineskip}%
  \addtolength\belowdisplayskip{-1\baselineskip}%
  
\begin{tikzpicture}
\node[] (data) {\includegraphics[width=0.7\textwidth]{figs/mauna_data.png}};
\node[below = 2cm of data] (post_param_helper) {};
\node[above = 1cm of post_param_helper] (post_param_helper_1) {};
\node[below = 1cm of post_param_helper] (post_param_helper_2) {};
\node[left = -2cm of post_param_helper] (post_param) {\includegraphics[width=0.6\textwidth]{figs/mauna_sample_1.png}};
\node[draw,rectangle,color=red,dashed,right = 2cm of post_param_helper,yshift=0.5cm] (zoom) {\includegraphics[width=0.2\textwidth]{figs/mauna_zoom.png}};
\node[draw,rectangle,color=red,dashed,left = 2.8cm of zoom, minimum width = 1.5cm, minimum height = 1.2cm] (zoom_in) {};
\node[below = 0.7cm of post_param_helper_2] (posterior) {\includegraphics[width=0.6\textwidth]{figs/mauna_structure.png}};

\node[draw,rectangle,below = 0.7cm of posterior] (formula) {\color{black}
$\mathbf{K}=\text{LIN} + \text{PER} + \text{SE} + \text{WN}$};
\node[draw,rectangle,below = 0.7cm of formula] (formula_param_1) {\color{black}\small
$= 2.7^2(x x^\prime) + 5.6^2 \exp \bigg( \frac{2 \sin^2 ( \pi (x - x^\prime)/3.7}{6.4^2} \bigg)
+ 0.4^2 \exp(-\frac{(x-x^\prime)^2}{2 \times 6.3^2}) +  1.9^2 \delta_{x,x^\prime} \label{eq:WN}$ };


\node[draw,rectangle,below = 0.7cm of formula_param_1,text width = \textwidth,minimum height = 1.5cm] (paragraph){\color{black}\footnotesize {\bf Qualitative Interpretation}:
The posterior peaks at a kernel structure with four additive components. Additive components hold globally, that is there are no higher level, qualitative aspects of the data that vary with the input space. The additive components are as follows: (i) a linearly increasing function or trend; (ii) a periodic function; (iii) a smooth function; and (iv) white noise.};







\node[draw, rectangle, left = -1.75cm of posterior,minimum width = 0.45cm, minimum height = 5.8cm,yshift=0.2cm] (mark_structure) {};
%\node[draw,very thick, rectangle, below = 1.1cm of data,minimum width = \textwidth, minimum height = 15cm] (posterior_frame) {};

%\node[left = 1.3cm of mark_structure] (paragraph_helper){};
\node[below =0.4cm of mark_structure,inner sep = 0pt,outer sep=0pt] (formula_helper) {};
\node[above =0.7cm of formula,inner sep = 0pt,outer sep=0pt] (formula_helper_2) {};

%\draw[-,dashed] (mark_structure.south) -- (formula_helper);
%\draw[-,dashed] (formula_helper) -- (formula_helper_2);
%\draw[->,dashed] (mark_structure) -- (paragraph_helper);
%\draw[->,dashed] (formula_helper_2) -- (formula);

\draw[->] (data) -- (post_param_helper_1);
\draw[->] (post_param_helper_2) -- (posterior);
\draw[->] (formula_helper_2) -- (formula);
\draw[-] (mark_structure) -- (formula_helper);
\draw[-] (formula_helper) -- (formula_helper_2);
\draw[->] (formula) -- (formula_param_1);
\draw[->] (formula_param_1) -- (paragraph);

\draw[->,dashed,red] (zoom_in) -- (zoom);

%\draw[->,line width=1pt,double distance=2pt] (data) -- (post_param);
\end{tikzpicture}
\addtolength\abovedisplayskip{1\baselineskip}%
\addtolength\belowdisplayskip{1\baselineskip}%



\caption{\small Structure Learning. Starting with raw data (a), we fit a \ac{GP}
(b) and compute the posterior distribution on structures (c). We take a sample
of the peak of this distribution ($\text{LIN}+\text{PER}+\text{SE}+\text{WN}$)
including its parameters and write it in functional form (d). We depict the
human readable interpretation (e). We used (d) to plot (b).}\label{fig:posterior}
\end{figure}
%%%%%%%%%%%%%%%%%%%%%%%%%%%%%%%%%%%%%%%%%%%%%%%%%%%%%%%%%%%%%%%%%%%%%%%%%%%%%%%%%
We illustrate results in Fig \ref{fig:posterior}. In Fig \ref{fig:posterior} (a) we depict the raw data. 
We see mean centered CO$_2$ measurements of the Mauna Loa Observatory, an atmospheric
baseline station on Mauna Loa, on the island of Hawaii. 
A description of the data set  can be found in  \citealp[][chapter 5]{rasmussen2006gaussian}.  
We use those raw data to compute a posterior on structure, parameters and \ac{GP}
samples.
The latter are shown in  Fig \ref{fig:posterior} (b)
where we zoom in to show how the posterior captures the error bars
adequately.
This posterior of the \ac{GP} is generated with a random sample from the parameters
of the peak of the distribution on structure (Fig \ref{fig:posterior} (c)).
We differentiate between a posterior distribution on kernel functions and on
distribution on symbolic expressions describing different kernel structures. 
This allows us to compute the posterior of symbollically equivalent structures,
such as $\Struct(\klin + \kper)=\Struct(\kper + \klin)$. Both structures yield and addition of a linear kernel and a periodic kernel, that is LIN + PER.
We coin the random variable over symbolic expressions for kernels as $\Ksrv$.
The distribution peaks at:
%%%%%%%%%%%%%%%%%%%%%%%%%%%%%%%%%%%%%%%%%%%%%%%%%%%%%%%%%%%%%%%%%%%%%%%%%%%%%%%%%
\begin{equation}
\Ksrv=\text{LIN} + \text{PER} + \text{SE} + \text{WN}.
\end{equation}
%%%%%%%%%%%%%%%%%%%%%%%%%%%%%%%%%%%%%%%%%%%%%%%%%%%%%%%%%%%%%%%%%%%%%%%%%%%%%%%%%
We write this kernel equation out in Fig \ref{fig:posterior} (d).
This kernel structure has a natural language interpretation that we spell out in
Fig \ref{fig:posterior} (e), explaining that 
the posterior peaks at a kernel structure with four additive components.
Each of which holds globally, that is there are no higher level, qualitative aspects
of the data that vary with the input space. The additive components for this result are as follows:
\begin{itemize}
\item a linearly increasing function or trend; 
\item a periodic function;
\item a smooth function; and
\item white noise.
\end{itemize}
 



Previous work on automated kernel discovery~\citep{duvenaud2013structure} illustrated the Mauna Loa data using an RQ kernel.
We resort to the white noise kernel instead of RQ (similar to \citep{lloyd2014automatic}).


%%%%%%%%%%%%%%%%%%%%%%%%%%%%%%%%%%%%%%%%%%%%%%%%%%%%%%%%%%%%%%%%%%%%%%%%%%%%%%%%%
%%%%%%%%%%%%%%%%%%%%%%%%%     Airline result   %%%%%%%%%%%%%%%%%%%%%%%%%%%%%%%%%%
%%%%%%%%%%%%%%%%%%%%%%%%%%%%%%%%%%%%%%%%%%%%%%%%%%%%%%%%%%%%%%%%%%%%%%%%%%%%%%%%
\myparagraph{Airline Data}
The second data set (Fig. \ref{fig:posterior_airline}) we depict results for is  the airline 
data set describing monthly totals of international airline passengers (\citealp{box2011time}, according to \citealp{duvenaud2013structure}). 
%%%%%%%%%%%%%%%%%%%%%%%%%%%%%%%%%%%%%%%%%%%%%%%%%%%%%%%%%%%%%%%%%%%%%%%%%%%%%%%%
\begin{figure}
\centering
 \addtolength\abovedisplayskip{-1\baselineskip}%
  \addtolength\belowdisplayskip{-1\baselineskip}%
  
\begin{tikzpicture}
\node (datatitle) {\small Raw Data};
\node[below = -0.25cm of datatitle](data) {\includegraphics[width=0.65\textwidth]{figs/airline_data.png}};
\node[below= 1.2cm of data] (post_param) {\includegraphics[width=0.65\textwidth]{figs/airline_sample_28.png}};


\node[below = 1.2cm of post_param] (posterior) {\includegraphics[width=0.55\textwidth]{figs/airline_structure.png}};


\node[draw,rectangle,below = 1.2cm of posterior] (formula_param_1) {\small\color{black}
$= 7.47^2(x x^\prime) +
\Bigg(0.27^2 \exp(-\frac{(x-x^\prime)^2}{2 \times 4.63^2}) \times 
7.34^2 \exp \bigg( \frac{2 \sin^2 ( \pi (x - x^\prime)/4.4}{4.55^2} \bigg)\Bigg)
+ 2.93^2 \delta_{x,x^\prime} \label{eq:WN}$ };


\node[draw,rectangle,below = 1.2cm of formula_param_1,text width =0.9\textwidth,minimum height = 1.5cm, font=\footnotesize] (paragraph){ 
The posterior peaks at a kernel structure with three additive components. Additive components hold globally, that is there are no higher level, qualitative aspects of the data that vary with the input space. The additive components are as follows: (i) a linearly increasing function or trend; (ii) a approximate periodic function; and (iv) white noise.};







\node[draw, rectangle, left = -1.7cm of posterior,minimum width = 0.45cm, minimum height = 5.2cm,yshift=0.2cm] (mark_structure) {};
%\node[draw,very thick, rectangle, below = 1.1cm of data,minimum width = \textwidth, minimum height = 15cm] (posterior_frame) {};

%\node[left = 1.3cm of mark_structure] (paragraph_helper){};
\node[below =0.45cm of mark_structure,inner sep = 0pt,outer sep=0pt] (formula_helper) {};
\node[above =1.2cm of formula_param_1,inner sep = 0pt,outer sep=0pt] (formula_helper_2) {};

%\draw[-,dashed] (mark_structure.south) -- (formula_helper);
%\draw[-,dashed] (formula_helper) -- (formula_helper_2);
%\draw[->,dashed] (mark_structure) -- (paragraph_helper);
%\draw[->,dashed] (formula_helper_2) -- (formula);

\draw[->] (data) -- node[right]{\small $\hat \fbf \sim
\mathcal{N}(\hat{\bm{\mu}},\hat\Kbf)$} (post_param_helper_1);
\draw[->] (post_param_helper_2) -- node[right]{\small Marginal Structure} (posterior);
\draw[->] (formula_helper_2) -- node[right] {\small
$\bm{\theta}=\{7.47,0.27,4.63,7.34,4.4,4.55,2.93\}$} (formula_param_1);
\draw[-] (mark_structure) -- node[left, yshift=-0.3cm] {\small $\Ktheta$} (formula_helper);
\draw[-] (formula_helper) --(formula_helper_2);
\draw[->] (formula_param_1) -- node[right]{\small Qualitative Interpretation} (paragraph);


\node[left=0.3cm of paragraph] (e){(e)}; 

\node[above=1.9cm of e] (d) {(d)}; 
\node[above=5.0cm of d] (c) {(c)}; 
\node[above=4.7cm of c] (b) {(b)}; 
\node[above=4.0cm of b] (a) {(a)}; 
\end{tikzpicture}
\addtolength\abovedisplayskip{1\baselineskip}%
\addtolength\belowdisplayskip{1\baselineskip}%



\caption{\small Structure Learning. Starting with raw data (a), we fit a \ac{GP}
(b) and compute the posterior distribution on structures (c). We take a sample
of the peak of this distribution ($\text{LIN}+\text{PER} \times \text{SE}+\text{WN}$)
including its parameters and write it in functional form (d). We depict the
human readable interpretation (e). We used (d) to plot (b).}\label{fig:posterior_airline}
\end{figure}
%%%%%%%%%%%%%%%%%%%%%%%%%%%%%%%%%%%%%%%%%%%%%%%%%%%%%%%%%%%%%%%%%%%%%%%%%%%%%%%%

We illustrate results for this data set in Fig \ref{fig:posterior_airline}. In Fig \ref{fig:posterior_airline} (a) we depict the raw data. 
Again, the data is mean centered and we use it to 
compute a posterior on structure, parameters and \ac{GP}
samples.
The latter are shown in  Fig \ref{fig:posterior_airline} (b).
This posterior of the \ac{GP} is generated with a random sample from the parameters
of the peak of the distribution on structure (Fig \ref{fig:posterior_airline} (c)).
The posterior over symbolic kernel expressions peaks at:
%%%%%%%%%%%%%%%%%%%%%%%%%%%%%%%%%%%%%%%%%%%%%%%%%%%%%%%%%%%%%%%%%%%%%%%%%%%%%%%%%
\begin{equation}
\Ksrv=\text{LIN} +  \text{SE} \times \text{PER}+ \text{WN}.
\end{equation}
%%%%%%%%%%%%%%%%%%%%%%%%%%%%%%%%%%%%%%%%%%%%%%%%%%%%%%%%%%%%%%%%%%%%%%%%%%%%%%%%%
We write this Kernel equation out in Fig \ref{fig:posterior_airline} (d).
This kernel structure has a natural language interpretation that we spell out in
Fig \ref{fig:posterior_airline} (e), explaining that 
the posterior peaks at a kernel structure with three additive components.
Additive components hold globally, that is there are no higher level, qualitative aspects
of the data that vary with the input space.
The additive components are as follows: 
\begin{itemize}
\item a linearly increasing function or trend;
\item a approximate periodic function; and
\item  white noise.
\end{itemize}
Both datasets served as illustrations in the Automatic Statistician project.



%%%%%%%%%%%%%%%%%%%%%%%%%%%%%%%%%%%%%%%%%%%%%%%%%%%%%%%%%%%%%%%%%%%%%%%%%%%%%%%%%
%%%%%%%%%%%%%%%%%%%%%%%%% Queries for time series %%%%%%%%%%%%%%%%%%%%%%%%%%%%%%%
%%%%%%%%%%%%%%%%%%%%%%%%%%%%%%%%%%%%%%%%%%%%%%%%%%%%%%%%%%%%%%%%%%%%%%%%%%%%%%%%%
%%%%%%%%%%%%%%%%%%%%%%%%%%%%%%%%%%%%%%%%%%%%%%%%%%%%%%%%%%%%%%%%%%%%%%%%%%%%%%%%%
\myparagraph{Querying time series}
With our Bayesian approach to structure learning we can gain valuable insights
into time series data that were previously unavailable.
This is due to our ability to estimate posterior marginal probabilities over the kernel structure.
Over this marginal, we define boolean search operations that allow us to query the data
for the probability of certain structures to hold true globally.
%%%%%%%%%%%%%%%%%%%%%%%%%%%%%%%%%%%%%%%%%%%%%%%%%%%%%%%%%%%%%%%%%%%%%%%%%%%%%%%%%
\begin{equation}
\label{eq:bool_present}
P(\Ksrv \mid \xbf,\ybf,\thetabf) = \frac{1}{T}
\sum\limits_{t=1}^T f(\Ksrv^t)\;\;\text{where}\, f(\Ksrv^t) = \begin{cases}
  1, & \text{if } \Ksrv \underset{global}{\in} \Ksrv^t, \\
  0, & \text{otherwise}.
\end{cases} 
\end{equation}
%%%%%%%%%%%%%%%%%%%%%%%%%%%%%%%%%%%%%%%%%%%%%%%%%%%%%%%%%%%%%%%%%%%%%%%%%%%%%%%%%
to ask whether it is true that a global structure $\Ksrv$ is present. $T$
is the number of all posterior samples for $\Ksrv$ and $\Ksrv^t$ is one such
sample. 
We can now ask simple questions, for example:
\begin{quotation}
Is there white noise in the data?
\end{quotation}
where we set $\Ksrv = $WN in (\ref{eq:bool_present}).
We can also formulate more sophisticated search operations using Boolean operators such as AND ($\land$) and OR ($\lor$).
The AND operator is defined as follows:
%%%%%%%%%%%%%%%%%%%%%%%%%%%%%%%%%%%%%%%%%%%%%%%%%%%%%%%%%%%%%%%%%%%%%%%%%%%%%%%%%
\[
P(\Ksrv^a \land \Ksrv^b \mid \xbf,\ybf, \thetabf)  = \frac{1}{N}
\sum\limits_{n=1}^N f(\Ksrv^t)\;\;
\]
%%%%%%%%%%%%%%%%%%%%%%%%%%%%%%%%%%%%%%%%%%%%%%%%%%%%%%%%%%%%%%%%%%%%%%%%%%%%%%%%%
where
\[
f(\Ksrv^t) = \begin{cases}
  1, & \text{if } \Ksrv^a\, \text{and}\, \Ksrv^b  \underset{global}{\in} \Ksrv^t, \\
  0, & \text{otherwise}\end{cases}.
\]
By estimating $P(\text{LIN} \land \text{WN} \mid \xbf, \ybf, \thetabf)$ we can use this operator to ask questions such as 
\begin{quotation}
Is there a Linear component AND a white noise in the data? 
\end{quotation}
Finally, we define the logical OR as
%%%%%%%%%%%%%%%%%%%%%%%%%%%%%%%%%%%%%%%%%%%%%%%%%%%%%%%%%%%%%%%%%%%%%%%%%%%%%%%%%
\begin{align*}
P(\Ksrv^a \lor \Ksrv^b \mid \xbf, \ybf, \thetabf)
=& P(\Ksrv^a \mid \xbf, \ybf, \thetabf) + P(\Ksrv^b \mid \xbf, \ybf, \thetabf)\\
 &- P(\Ksrv^a \land \Ksrv^b \mid \xbf, \ybf, \thetabf)
\end{align*}
%%%%%%%%%%%%%%%%%%%%%%%%%%%%%%%%%%%%%%%%%%%%%%%%%%%%%%%%%%%%%%%%%%%%%%%%%%%%%%%%%
which allows us to ask questions about structures that are logically connected with OR, such as:
\begin{quotation}
Is there white noise or heteroskedastic noise?
\end{quotation}
by estimating $P(\text{LIN} \times \text{WN}\;\;{\large\lor}\;\; \text{WN} \mid
\xbf, \ybf, \thetabf)$.
We know that noise can either be heteroskedastic or white,
and we also know due to simple manipulations using kernel algebra
that  $\text{LIN} \times \text{WN}$ and $\text{WN}$ are the only possible ways to construct noise with kernel composition, we see that we can generalize the 
question above to:
\begin{quotation}
Is there noise in the data? 
\end{quotation}
where we write the marginal posterior on qualitative structure for noise:
%%%%%%%%%%%%%%%%%%%%%%%%%%%%%%%%%%%%%%%%%%%%%%%%%%%%%%%%%%%%%%%%%%%%%%%%%%%%%%%%%
\begin{equation}
P(\Ksrv^{\text{noise}} \mid \xbf, \ybf, \thetabf) = P(\text{LIN} \times
\text{WN}\;\;{\large\lor}\;\; \text{WN} \mid \xbf, \ybf, \thetabf).
\end{equation}
%%%%%%%%%%%%%%%%%%%%%%%%%%%%%%%%%%%%%%%%%%%%%%%%%%%%%%%%%%%%%%%%%%%%%%%%%%%%%%%%%
Note that this allows us to start with general queries and 
subsequently formulate follow up queries that go into more detail.
For example, we could start with a general query, such as:
\begin{quotation}
What is the probability of a trend, a recurring pattern {\bf and} noise in the data?
\end{quotation}
and then follow up with more detailed questions (Fig \ref{fig:query}).
%%%%%%%%%%%%%%%%%%%%%%%%%%%%%%%%%%%%%%%%%%%%%%%%%%%%%%%%%%%%%%%%%%%%%%%%%%%%%%%%%
\begin{figure}
\centering
\footnotesize
% \includegraphics[width=.153\textwidth]{figs/gpSamples/main.png}
\begin{tabular}{cccccc}
\multicolumn{6}{c}{ \includegraphics[width=0.6\textwidth]{figs/mauna_data.png}}     \\                                               
\multicolumn{6}{c}{\tikzmark{a}}     \\  
\multicolumn{6}{c}{}     \\   
\multicolumn{6}{c}{\tikzmark{b}}     \\  
\multicolumn{6}{c}{ \includegraphics[width=0.5\textwidth]{figs/mauna_structure.png}}     \\                                               
\multicolumn{6}{c}{\tikzmark{c}}     \\  
\multicolumn{6}{c}{}     \\   
\multicolumn{6}{c}{\tikzmark{d}}     \\           
\multicolumn{6}{c}{\normalsize \color{blue} What is the probability of a trend, a recurring pattern {\bf and} noise in the data?}     \\               
\multicolumn{6}{c}{$P\big((\text{LIN}\lor\text{LIN}\times\text{SE})\land
(\text{PER}\lor\text{PER}\times\text{SE}\lor\text{PER}\times\text{LIN})\land
(\text{WN}\lor\text{LIN}\times\text{WN})\big) = 0.36
$}                                                     \\
         \multicolumn{2}{c}{ \tikzmark{trend_part}}     &  \multicolumn{2}{c}{$\;\;\;\;$\tikzmark{recurring_part}}   &     \multicolumn{2}{c}{\tikzmark{noise_part}}                               \\
          &          &  &  &  &                            \\
\multicolumn{2}{c}{\tikzmark{trend}}& & &   \multicolumn{2}{c}{ \tikzmark{noise}} \\  
\multicolumn{2}{c}{\normalsize \color{blue} Is there a trend?}& & &   \multicolumn{2}{c}{ \normalsize \color{blue} Is there noise?} \\
\multicolumn{2}{c}{$P(\text{LIN}\lor\text{LIN}\times\text{SE}) = 0.65$}& & &   \multicolumn{2}{c}{$P(\text{WN}\lor\text{LIN}\times\text{WN}) = 0.75$} \\
\multicolumn{2}{c}{\includegraphics[width=.15\textwidth]{figs/gpSamples/trend.png}}& & &   \multicolumn{2}{c}{\includegraphics[width=.15\textwidth]{figs/gpSamples/noise.png}} \\
\multicolumn{2}{c}{\tikzmark{trend_below}}& & &   \multicolumn{2}{c}{ \tikzmark{noise_below}} \\  
        &          &  &       &       &                            \\
         \tikzmark{lin}  &\tikzmark{linse}             &         &       &                 \tikzmark{linwn}  &\tikzmark{wn}            \\

           {\normalsize \color{blue} A linear trend?}  &{\normalsize \color{blue} A smooth trend?}   &         &       &                 {\normalsize \color{blue} Heteroskedastic noise?}  &{\normalsize \color{blue} White noise?}           \\
          $P(\text{LIN}) = 0.63$& $P(\text{LIN}\times\text{SE}) = 0.02$         &          &    &$P(\text{LIN}\times\text{WN}) = 0$  &$P(\text{WN}) = 0.75$    \\
                \includegraphics[width=.15\textwidth]{figs/gpSamples/lin.png}& \includegraphics[width=.15\textwidth]{figs/gpSamples/selin.png}    &          &    &\includegraphics[width=.15\textwidth]{figs/gpSamples/linwn.png}  &\includegraphics[width=.15\textwidth]{figs/gpSamples/wn.png}   \\
         &          &         &\tikzmark{recurring}       &                &          \\
      \multicolumn{6}{c}{\normalsize \color{blue} Is there repeating structure?} \\
  \multicolumn{6}{c}{$P(\text{PER}\lor\text{PER}\times\text{SE}\lor\text{PER}\times\text{LIN}) = 0.73$} \\
   \multicolumn{6}{c}{\includegraphics[width=.15\textwidth]{figs/gpSamples/recurring.png}$\;\;\;\;\;\;$} \\
     & &    \multicolumn{2}{c}{\tikzmark{recurring_below}} & &\\
      &          &  &       &       &                            \\
         &  \tikzmark{per}        &   \multicolumn{2}{c}{\tikzmark{perse}}     &   \tikzmark{perlin}    &    \\
        \multicolumn{2}{c}{$P(\text{PER}) = 0.32$}         &           \multicolumn{2}{c}{$P(\text{PER}\times\text{SE}) = 0.34$}         &                   \multicolumn{2}{c}{$P(\text{PER}\times\text{LIN}) = 0.07$}           \\
        \multicolumn{2}{c}{\includegraphics[width=.15\textwidth]{figs/gpSamples/per.png}}         &           \multicolumn{2}{c}{\includegraphics[width=.15\textwidth]{figs/gpSamples/seper.png}}         &                   \multicolumn{2}{c}{\includegraphics[width=.15\textwidth]{figs/gpSamples/perlin.png}}           
        \end{tabular}
 \begin{tikzpicture}[overlay, remember picture, yshift=.25\baselineskip, shorten >=.5pt, shorten <=.5pt]
   \draw [->] ({pic cs:a}) to ({pic cs:b});
      \draw [->] ({pic cs:c}) to ({pic cs:d});
    \draw [->] ({pic cs:recurring_part}) to ({pic cs:recurring});
    
     \draw [->] ({pic cs:recurring_below}) to ({pic cs:per});
      \draw [->] ({pic cs:recurring_below}) to ({pic cs:perse});
     \draw [->] ({pic cs:recurring_below}) to ({pic cs:perlin});
    
    
      \draw [->] ({pic cs:trend_part}) to ({pic cs:trend});
        \draw [->] ({pic cs:trend_below}) to ({pic cs:lin});
    \draw [->] ({pic cs:trend_below}) to ({pic cs:linse});
    
     \draw [->] ({pic cs:noise_part}) to ({pic cs:noise});
       \draw [->] ({pic cs:noise_below}) to ({pic cs:wn});
         \draw [->] ({pic cs:noise_below}) to ({pic cs:linwn});
  \end{tikzpicture}
%  \multicolumn{3}{c}{$P(\text{PER}\lor\text{PER}\times\text{SE}\lor\text{PER}\times\text{LIN})$}
\caption{{\bf Querying structural motifs in in time series using posterior inference
over kernel structure.} The kernel structure serves as a way to formulate
natural language questions about the data (blue). The initial question of interest
(top) is a fairly
general one: "What is the probability of a trend, a recurring
pattern and noise in the data?" Below the natural language version of this
question, the same question is formulated as an inference problem (black) over the
marginal probability on kernels with Boolean operators AND ($\land$) and OR ($\lor$). 
To gain  a deeper understanding of specific motifs in the time series more specific queries can
be written.
On the right, a query asks whether there is noise in the data (blue) by computing the disjunction of the marginal
of a global white noise kernel and a multiplication between a linear and a white
noise kernel (black). Samples from the predictive prior $\ybf_*$ of such kernels give an
indication of the qualitative aspects that a kernel structure implies (coloured curves below
the marginal). 
If the probability that there is noise in the data is high then it makes sense
to drill even deeper asking more detailed questions. With regards to noise, this
translates to querying whether or not the data supports the hypothesis that there is
heteroskedastic noise or white noise. Queries for motifs of repeating structure
are shown in the middle of the tree, queries related to trends on the left.}\label{fig:query}
\end{figure}
%%%%%%%%%%%%%%%%%%%%%%%%%%%%%%%%%%%%%%%%%%%%%%%%%%%%%%%%%%%%%%%%%%%%%%%%%%%%%%%%%
This way of querying your data for their statistical implications is in stark contrast to what previous research in automatic kernel construction was able to provide.
We could view our approach as a time series search engine which allows us to test whether or not certain structures can be found
in an available time series.
Another way to view this approach is as a new language to interact with the world.
Real-world observations often come with time-stamps and in form
of continuous valued sensor measurements.  
We provide the toolbox to query such observations in a similar manner as
one would query a knowledge base in a logic programming language.








\FloatBarrier
\section{Bayesian Optimization}
Bayesian optimization casts the problem of finding the global maximum of an unknown function as a hierarchical decision problem~\citep{ghahramani2015probabilistic}.
Evaluating the actual function may be very expensive, either in computation time or in some other resource.
For one example, when searching for the best configuration for the learning algorithm of a large convolutional neural network, a large amount of computational work is required to evaluate a candidate configuration, and the space of possible configurations is high-dimensional.
Another common example, alluded to in Section \ref{sec:gpmem-broader}, is data acquisition: for machine learning problems in which a large body of data is available, it is often desirable to choose the right queries to produce a data set on which learning will be most effective.
In continuous settings, many Bayesian optimization methods employ GPs~\citep[e.g.][]{snoek2012practical}.

To demonstrate Bayesian optimization with \gpmem\, we have implemented a version of Thompson sampling using GPs in Venture.
Thompson sampling~\citep{thompson1933likelihood} is a widely-used Bayesian framework for solving exploration-exploitation problems.
Our implementation has two notable features: (i) the ability to search over a broader space of contexts than the parametric families that are typically used, and (ii) the parsimony of the resulting probabilistic program.


\subsection{Thompson sampling framework}
We now lay out the setup of Thompson sampling for Markov decision processes (MDPs).
An agent is to take a sequence of actions $a_1, a_2, \ldots$ from a (possibly infinite) set of possible actions $\Acal$.
After each action, a reward $r \in \R$ is received, according to an unknown conditional distribution $P_\true(r|a)$.
The agent's goal is to maximize the total reward received for all actions.
In Thompson sampling, the Bayesian agent accomplishes this by placing a prior distribution $P(\theta)$ on the possible ``contexts'' $\theta \in \Theta$.
Here a context is a believed model of the conditional distributions $\{P(r|a)\}_{a \in \Acal}$, or at least, a believed statistic of these conditional distributions which is sufficient for deciding an action $a$.
One example of such a sufficient statistic is the conditional mean $V(a|\theta) = \Ebkt{r|a,\theta}$, which can be thought of as a value function.
Thompson sampling thus has the following steps, repeated as long as desired:
\begin{enumerate}
  \item Sample a context $\theta \sim P(\theta)$.
  \item Choose an action $a \in \Acal$ which (approximately) maximizes $V(a|\theta) = \Ebkt{r|a,\theta}$.
  \item\label{itm:Thompson-conditioning}
    Let $r_\true$ be the reward received for action $a$.
    Update the believed distribution on $\theta$, i.e., $P(\theta) \gets P_\rmnew(\theta)$ where $P_\rmnew(\theta) = P\pn{\theta \mvert a \mapsto r_\true}$.
\end{enumerate}
Note that when $\Ebkt{r|a,\theta}$ (under the sampled value of $\theta$ for some points $a$) is far from the true value $\Ebkt[P_\true]{r|a}$, the chosen action $a$ may be far from optimal, but the information gained by probing action $a$ will improve the belief $\theta$.
This amounts to ``exploration.''
When $\Ebkt{r|a,\theta}$ is close to the true value except at points $a$ for which $\Ebkt{r|a,\theta}$ is low, exploration will be less likely to occur, but the chosen actions $a$ will tend to receive high rewards.
This amounts to ``exploitation.''
Roughly speaking, exploration will happen until the context $\theta$ is reasonably sure that the unexplored actions are probably not optimal, at which time the sampler will exploit by choosing actions in regions it knows to have high value.

Typically, when Thompson sampling is implemented, the search over contexts $\theta \in \Theta$ is limited by the choice of representation.
In traditional programming environments, $\theta$ often consists of a few numerical parameters for a family of distributions of a fixed functional form.
With work, a mixture of a few functional forms is possible; but without probabilistic programming machinery, implementing a rich context space $\Theta$ would be an unworkably large technical burden.
In a probabilstic programing language, however, the representation of heterogeneously structured or infinite-dimensional context spaces is quite natural.
Any computable model of the conditional distributions $\br{P(r|a)}_{a \in \Acal}$ can be represented as a stochastic procedure $(\lambda (a) \ldots)$.
Thus, for computational Thompson sampling, the most general context space $\widehat\Theta$ is the space of program texts.
Any other context space $\Theta$ has a natural embedding as a subset of $\widehat\Theta$.


\subsection{Thompson sampling in Venture}

Because Venture supports sampling and inference on (stochastic-)procedure-valued random variables (and the generative models which produce those procedures), Venture can capture arbitrary context spaces as described above.
To demonstrate, we have implemented Thompson sampling in Venture in which the contexts $\theta$ are Gaussian processes over the action space $\Acal = \R$.
That is, $\theta = (\mu, K)$, where the mean $\mu$ is a computable function $\Acal \to \R$ and the covariance $K$ is a computable (symmetric, positive-semidefinite) function $\Acal \times \Acal \to \R$.
This represents a Gaussian process $\br{R_a}_{a \in \Acal}$, where $R_a$ represents the reward for action $a$.
% where for any finite subset $\br{a_i}_{i=1}^{n} \subset \Acal$, the marginal distribution on $\pn{R_{a_i}}$ is Gaussian with mean $\pn{\mu(a_i)}_{i=1}^{n}$ and covariance matrix $\begin{pmatrix} K(a_i,a_j) \end{pmatrix}_{1 \leq i,j \leq n}$.
Computationally, we represent a context not as a pair of infinite lookup tables for $\mu$ and $K$, but as a finite data structure $\theta = (K_\prior, \sigma, \ell, \abf_\past, \rbf_\past)$, where
\begin{itemize}
  \item $K_\prior = K_{\prior,\sigma,\ell}$ is a procedure, with parameters $\sigma,\ell$, to be used as the prior covariance function: $K_\prior(a,a') = \sigma^2 \exp\pn{-\frac{(a-a')^2}{2\ell^2}}$
  \item $\sigma$ and $\ell$ are (hyper)parameters for $K_\prior$
  \item $\abf_\past = \pn{a_i}_{i=1}^{n}$ are the previously probed actions
  \item $\rbf_\past = \pn{r_i}_{i=1}^{n}$ are the corresponding rewards
\end{itemize}
To simplify the treatment, we take prior mean $\mu_\prior \equiv 0$.  The mean and covariance for $\theta$ are then gotten by the usual conditioning formula:
\begin{align*}
  \mu(\abf)
  &= \mu\pn{\abf \mvert \abf_\past, \rbf_\past} \\
  &= K_\prior(\abf, \abf_\past)
     \,K_\prior(\abf_\past, \abf_\past)^{-1}
     \,\rbf_\past \\
  K(\abf, \abf)
  &= K\pn{\abf, \abf \mvert \abf_\past, \rbf_\past} \\
  &= K_\prior(\abf, \abf)
     - K_\prior(\abf, \abf_\past)
       \,K_\prior(\abf_\past, \abf_\past)^{-1}
       \,K_\prior(\abf_\past, \abf).
\end{align*}
Note that even in this simple example, the context space $\Theta$ is not a finite-dimensional parametric family, since the vectors $\abf_\past$ and $\rbf_\past$ grow as more samples are taken.
$\Theta$ is, however, quite easily representable as a computational procedure together with parameters and past samples, as we do in the representation $\theta = (K_\prior, \sigma, \ell, \abf_\past, \rbf_\past)$.


\subsection{Implementation with \gpmem}

As a demonstration, we use Thompson sampling to optimize an unknown function $V(x)$ (the value function) using \gpmem.
(\textbf{TODO} we should not assume $V$ is deterministic, it would be easy enough to make it random or have it give noisy samples.)
We assume $V$ is made available to Venture as a black-box.
The code for optimizing $V$ is given in Listing \ref{alg:bayesopt}.
For step \ref{itm:Thompson-conditioning} of Thompson sampling, the Bayesian update, we not only condition on the new data (the chosen action $a$ and the received reward $r$), but also perform inference on the hyperparameters $\sigma, \ell$ using a Metropolis--Hastings sampler.
These two inference steps take 1 line of code: 0 lines to condition on the new data (as this is done automatically by \gpmem), and 1 line to call Venture's built-in \texttt{MH} operator.
The results are shown in Figure \ref{fig:bayesopt-sequence}.
We can see from the figure that, roughly speaking, each successive probe point $a$ is chosen either because the current model $V_\emu$ thinks it will have a high reward, or because the value of $V_\emu(a)$ has high uncertainty.
In the latter case, probing at $a$ decreases this uncertainty and, due to the smoothing kernel, also decreases the uncertainty at points near $a$.
We thus see that our Thompson sampler simultaneously learns the value function and optimizes it.

\begin{minipage}{\linewidth}
\small
\belowcaptionskip=-10pt
\begin{lstlisting}[frame=single,caption={
  Code for Bayesian optimization using \gpmem.
  %The procedure \texttt{V\_compute} probes \texttt{V} directly, thus improving the GP model \texttt{V\_emu}.
  %(\texttt{V\_emu\_pointwise} is simply a shortcut for sampling the GP model at a single point; \texttt{V\_emu} is more general, allowing joint samples to be taken at any set of points.)
  In the loop, \texttt{V\_compute} is called to probe the value of \texttt{V} at a new argument.
  The new argument, \texttt{(mc\_argmax V\_emu\_pointwise mc\_sampler)}, is a Monte Carlo estimate of the maximum pointwise sample of \texttt{V\_emu} (itself a stochastic quantity), with the Monte Carlo samples being drawn in this case uniformly between $-20$ and $20$.
  After each new call to \texttt{V\_compute}, the Metropolis--Hastings algorithm is used to perform inference on the hyperparameters of the covariance function in the GP model in light of the new conditioning data.
  Once enough calls to \texttt{V\_compute} have been made (in our case we stopped at 15 calls), we can inspect the full list of probed $(a,r)$ pairs with \texttt{extract\_stats}.
  The answer to our maximization problem is simply the pair having the highest $r$; but our algorithm also learns more potentially useful information.},mathescape,numbers=none,label=alg:bayesopt]
assume sf = tag(quote(hyper), 0, uniform_continuous(0, 10))
assume l = tag(quote(hyper), 1, uniform_continuous(0, 10))
assume se = make_squaredexp(sf, l)
assume blackbox_f = get_bayesopt_blackbox()
assume (f_compute, f_emu) = gpmem(blackbox_f, se)

define get_uniform_candidate = proc(prev_xs) {
  uniform_continuous(-20, 20)
}

define mc_argmax = proc(func, prev_xs) {
  // Monte Carlo estimator for the argmax of func.
  run(do(
    candidate_xs <- mapv(proc(i) {get_uniform_candidate(prev_xs)},
                         linspace(0, 19, 20)),
    candidate_ys <- mapv(func, candidate_xs),
    lookup(candidate_xs, argmax_of_array(candidate_ys))))
}

define emulator_point_sample = proc(x) {
  run(sample(lookup(
    f_emu(array(x)),
    0)))
}

infer repeat(15, do(pass,
     // Phase 2: Call f_compute on the next probe point
     predict f_compute(
                 mc_argmax(emulator_point_sample, '_)),
     // Phase 1: Hyperparameter inference
     mh(quote(hyper), one, 50)))
\end{lstlisting}
\end{minipage}
      
\begin{figure}
\centering
    \includegraphics[height=0.8\textheight]{figs/BayesOpt_gpmem_sequence.png}
    \caption{
      Dynamics of Thompson sampling in Venture.
      The blue curve is the true function $V$, and the red region is a blending of 100 samples of the curve generated (jointly) by a GP-based emulator $V_\emu$.
      The left and right columns show the state of $V_\emu$ before and after hyperparameter inference is run on the new data, respectively.
      (We can see, for example, that after the seventh probe point, the Metropolis--Hastings sampler chose a ``crazy'' set of hyperparameters, which was corrected at the next inference step.)
      In the right column, the next chosen probe point is circled in green.
      Each successive probe point $a$ is the (stochastic) maximum of $V_\emu$, sampled pointwise and conditioned on the values of the previously probed points.
      Note that probes tend to happen at points either where the value of $V_\emu$ is high, or where $V_\emu$ has high uncertainty.
      }
    \label{fig:bayesopt-sequence}
\end{figure}



%We consider a true and  unknown reward function $r(x)$ that we estimate with a GP prior $\mathcal{GP}(0,K(\mathbf{x},\mathbf{x}))$. We denote past observations with $\mathcal{D} = \{(x;

\FloatBarrier

\section{Conclusion}
This paper has shown that it is feasible and useful to embed Gaussian processes
in higher-order probabilistic programming languages by treating them as a kind
of statistical memoizer. It has described classic GP regression with both fully
Bayesian and MAP inference in a hierarchical hyperprior, as well as state-of-the-art
applications to discovering symbolic structure in time series and to Bayesian optimization.
All the applications share a common 100-line Python GP library and require fewer than 20 lines
of probabilistic code each.

These results suggest several research directions. First, it will be important
to develop versions of {\tt gpmem} that are optimized for larger-scale
applications. Possible approaches include the standard low-rank approximations
to the kernel matrix that are popular in machine learning~\citep{bui2014tree} as well
as more sophisticated sampling algorithms for approximate conditioning of the
GP~\citep{lawrence2009efficient}.
Second, it seems fruitful to abstract the notion of a ``generalizing" memoizer
from the specific choice of a Gaussian process model as the mechanism for
generalization. ``Generalizing" or statistical memoizers with custom regression techniques could be broadly useful in performance engineering and scheduling systems.
The timing data from performance benchmarks could be run through a generalizing memoizer by default.
This memoizer could be queried (and its output error bars examined) to inform the best strategy
for performing the computation or predict the likely runtime of long-running jobs.
 Third, the structure learning application suggests follow-on research in information
retrieval for structured data. It should be possible to build a time series search engine
that can handle search predicates such as ``has a rising trend starting around
1988" or ``is perodic during the 1990s".
The variation on the Automated Statistician presented in this paper can provide ranked result
sets for these sorts of queries because it tracks posterior uncertainty over structure and also
because the space of structural patterns that it can handle is easy to modify by making small
changes to a short VentureScript program.


The field of Bayesian nonparametrics offers a principled, fully Bayesian
response to the empirical modeling philosophy in machine learning~\citep{ghahramani2013bayesian},
where Bayesian inference is used to encode a state of broad ignorance rather
than a bias stemming from strong prior knowledge. It is perhaps surprising that
two key objects from Bayesian nonparametrics, Dirichlet processes and \ac{GP}s,
fit naturally in probabilistic programming as variants of
memoization~\citep{roy2008stochastic}. It is not yet clear if the same will be true
for other processes, e.g. Wishart processes, or hierarchical Beta processes. We hope that the results in this paper encourage the development of other nonparametric libraries for higher-order probabilistic programming languages.


\section*{Appendix}
\subsection*{A Covariance Functions}
SE and WN are defined in the text above, for completeness we will introduce the covariance:
% ToDo: check if this should be removed
\begin{equation}
k_{LIN}(x,x^\prime) = \theta (x x^\prime) 
\end{equation}
\begin{equation}
k_{PER}(x,x^\prime) = \theta \exp \bigg( \frac{2 \sin^2 ( \pi (x - x^\prime)/p}{\ell^2} \bigg) 
\end{equation}
\begin{equation}
k_{RQ}(x,x^\prime) =   \theta \bigg(1 + \frac{(x - x^\prime)^2}{2 \alpha \ell^2} \bigg)^{-\alpha}
\end{equation}


\subsection*{B Covariance Simplification}
\begin{minipage}{\linewidth}
\small
\belowcaptionskip=-10pt
\begin{lstlisting}[frame=single,mathescape,label=alg:simplify,basicstyle=\selectfont\ttfamily]

SE $\times$ SE                  $\rightarrow$ SE 
{SE,PER,C,WN} $\times$ WN       $\rightarrow$ WN
LIN $+$ LIN                $\rightarrow$ LIN
{SE,PER,C,WN,LIN} $\times$ C    $\rightarrow$  {SE,PER,C,WN,LIN} 
\end{lstlisting}
\end{minipage}
Rule 1 is derived as follows:
\begin{equation}
\begin{aligned}
\sigma_c^2 \exp(-\frac{(x-x^\prime)^2}{2\ell_c^2})  &=  \sigma_a^2 \exp(-\frac{(x-x^\prime)^2}{2\ell_a^2}) \times  \sigma_b^2 \exp(-\frac{(x-x^\prime)^2}{2\ell_b^2}) \\
&= \sigma_c^2 \exp(-\frac{(x-x^\prime)^2}{2\ell_a^2}) \times   \exp(-\frac{(x-x^\prime)^2}{2\ell_b^2}) \\
&= \sigma_c^2 \exp \bigg(-\frac{(x-x^\prime)^2}{2\ell_a^2} -\frac{(x-x^\prime)^2}{2\ell_b^2}\bigg) \\
&= \sigma_c^2 \exp \bigg(-\frac{(x-x^\prime)^2}{2\ell_c^2}\bigg) \\
\end{aligned}
\end{equation}
Rule 3 is derived as follows:
\begin{equation}
 \theta_c (x \times x^\prime) = \theta_a (x \times x^\prime) + \theta_b (x \times x^\prime) 
\end{equation}
For stationary kernels that only depend on the lag vector between $x$ and $x^\prime$ it holds that multiplying such a kernel with a WN kernel we get another WN kernel. Take for example the SE kernel:
\begin{equation}
 \sigma_a^2 \exp \bigg(-\frac{(x-x^\prime)^2}{2\ell_c^2}\bigg) \times  \sigma_b \delta_{x,x^\prime} =  \sigma_a \sigma_b \delta_{x,x^\prime}
\end{equation}
Multiplying any kernel with a constant obviously changes only the scale parameter of a kernel.


\newpage
\bibliography{May2015}
\bibliographystyle{apalike}
\end{document}
